\documentclass[11pt, a4paper]{article}

% =====================================================================
% PACCHETTI FONDAMENTALI (pienamente compatibili con PdfLaTeX su MiKTeX)
% =====================================================================
\usepackage[utf8]{inputenc}
\usepackage[T1]{fontenc}
\usepackage[italian]{babel}
\usepackage{lmodern}
\usepackage{microtype}
\usepackage[a4paper, top=2.5cm, bottom=2.5cm, left=2.5cm, right=2.5cm, headheight=14pt]{geometry}

% Impaginazione, colori, tabelle, listati di codice, titoli ed elenchi
\usepackage{xcolor}
\usepackage{titlesec}
\usepackage{fancyhdr}
\usepackage{lastpage}
\usepackage{enumitem}
\usepackage{parskip}
\usepackage{graphicx}
\usepackage{array}
\usepackage{tabularx}
\usepackage{longtable}
\usepackage{listings}

% Collegamenti cliccabili nell'indice e metadati PDF (da caricare per ultimo)
\usepackage[
	hidelinks,
	pdfauthor={Alberto Mancini},
	pdftitle={Analisi e Superamento delle Tecniche Anti-VM e Anti-Sandbox nei Malware PlugX, RAT in Go Generico, Astaroth e Latrodectus}
]{hyperref}

% =====================================================================
% PALETTE COLORI, STILE DEI TITOLI E CONFIGURAZIONE LISTATI
% =====================================================================
\definecolor{primary}{HTML}{1B365D}    % Blu accademico per titoli e dettagli
\definecolor{muted}{HTML}{5A6675}      % Grigio per intestazioni e metadati
\definecolor{boxbg}{HTML}{F3F6FA}      % Sfondo leggero per i riquadri di output
\definecolor{codebg}{HTML}{F8FAFC}     % Sfondo chiaro per i listati di codice
\definecolor{rulecol}{HTML}{CBD5E1}    % Grigio chiaro per le linee divisorie e bordi
\definecolor{kwcol}{HTML}{0550AE}      % Blu intenso per keyword Python
\definecolor{strcol}{HTML}{0A3069}     % Blu scuro per stringhe
\definecolor{cmtcol}{HTML}{57606A}     % Grigio ardesia per i commenti

\titleformat{\section}
{\Large\bfseries\color{primary}}
{\thesection.}{0.6em}{}
[\vspace{2pt}\color{rulecol}\titlerule]

\titleformat{\subsection}
{\large\bfseries\color{primary!85!black}}
{\thesubsection}{0.6em}{}

\titleformat{\subsubsection}
{\normalsize\bfseries\color{black!85}}
{\thesubsubsection}{0.6em}{}

% Spaziatura ordinata per elenchi, paragrafi tecnici e tabelle
\setlength{\emergencystretch}{2.5em}
\setlist[itemize]{leftmargin=1.5em, itemsep=3pt, topsep=4pt}
\setlist[itemize,2]{leftmargin=1.5em, itemsep=2pt, topsep=3pt, before=\raggedright}
\setlist[enumerate]{leftmargin=1.8em, itemsep=4pt, topsep=4pt}
\renewcommand{\arraystretch}{1.28}
\renewcommand{\lstlistingname}{Listato}

% Colonne flessibili allineate a sinistra per evitare spaziature anomale e sforamenti nelle tabelle
\newcolumntype{Y}{>{\raggedright\arraybackslash}X}
\newcolumntype{P}[1]{>{\raggedright\arraybackslash}p{#1}}
\newcolumntype{C}[1]{>{\centering\arraybackslash}p{#1}}

% Stile per i listati Python degli harness Qiling
\lstdefinestyle{pythoncode}{
	language=Python,
	backgroundcolor=\color{codebg},
	basicstyle=\footnotesize\ttfamily\color{black!90},
	keywordstyle=\bfseries\color{kwcol},
	stringstyle=\color{strcol},
	commentstyle=\itshape\color{cmtcol},
	morekeywords={True,False,None,as,with,yield,async,await,self},
	emph={winsdkapi,Qiling,STDCALL,CDECL,DWORD,HANDLE,POINTER,SIZE_T,ULONG,ULONGLONG,BOOL,INT,UINT,WORD,BYTE,STRING,WSTRING,LPCSTR,LPVOID},
	emphstyle=\bfseries\color{primary!80!black},
	numbers=left,
	numberstyle=\tiny\color{muted},
	numbersep=8pt,
	frame=single,
	rulecolor=\color{rulecol},
	framesep=5pt,
	xleftmargin=14pt,
	xrightmargin=4pt,
	aboveskip=8pt,
	belowskip=8pt,
	breaklines=true,
	breakatwhitespace=false,
	columns=fullflexible,
	keepspaces=true,
	showstringspaces=false,
	tabsize=4,
	captionpos=b
}

% Stile per gli output di console, log di esecuzione e report JSON/telemetrici
\lstdefinestyle{consoleout}{
	backgroundcolor=\color{boxbg},
	basicstyle=\scriptsize\ttfamily\color{black!90},
	frame=leftline,
	framerule=2.2pt,
	rulecolor=\color{primary!75},
	framesep=6pt,
	xleftmargin=8pt,
	xrightmargin=4pt,
	aboveskip=7pt,
	belowskip=7pt,
	breaklines=true,
	breakatwhitespace=false,
	columns=fullflexible,
	keepspaces=true,
	showstringspaces=false,
	tabsize=2,
	captionpos=b
}

% =====================================================================
% INTESTAZIONE E PIE' DI PAGINA
% =====================================================================
\pagestyle{fancy}
\fancyhf{}
\renewcommand{\headrulewidth}{0.4pt}
\renewcommand{\footrulewidth}{0.4pt}
\fancyhead[L]{\small\color{muted}Relazione Tecnica -- Analisi Malware Anti-VM e Anti-Sandbox}
\fancyhead[R]{\small\color{muted}Alberto Mancini}
\fancyfoot[L]{\small\color{muted}Matricola: 352637}
\fancyfoot[R]{\small\color{muted}Pagina \thepage\ di \pageref{LastPage}}

% =====================================================================
% INIZIO DOCUMENTO
% =====================================================================
\begin{document}

\pagenumbering{gobble}
\begin{titlepage}
	\centering
	
	% --- 1. INTESTAZIONE UNIVERSITA' E CORSO ---
	{\Large \textsc{Universit\`a degli Studi di Perugia}} \\[0.3cm]
	{\large Dipartimento di Matematica e Informatica} \\[0.2cm]
	{\large Corso di Laurea Triennale in Informatica} \\[1.2cm]
	
	% --- 2. LOGO ---
	\includegraphics[width=0.42\textwidth]{img/unipg_icon.png}
	
	% --- 3. TIPO DI DOCUMENTO E TITOLO ---
	\vfill
	
	{\large \textbf{Relazione di Tirocinio}} \\[0.4cm]
	\rule{\textwidth}{0.5pt} \\[0.4cm]
	{\LARGE \bfseries Analisi e Superamento delle Tecniche Anti-VM e Anti-Sandbox nei Malware PlugX, RAT in Go Generico, Astaroth e Latrodectus \par}
	\vspace{0.3cm}
	\rule{\textwidth}{0.5pt}
	
	\vfill
	
	% --- 4. RELATORE (SINISTRA) E STUDENTE (DESTRA) ---
	\begin{minipage}[t]{0.48\textwidth}
		\begin{flushleft}
			\textbf{Tutor:} \\
			Prof. Francesco Santini
		\end{flushleft}
	\end{minipage}
	\hfill
	\begin{minipage}[t]{0.48\textwidth}
		\begin{flushright}
			\textbf{Studente:} \\
			Alberto Mancini \\
			Matricola: 352637
		\end{flushright}
	\end{minipage}
	
	\vspace{1.8cm}
	
	% --- 5. ANNO ACCADEMICO ---
	\rule{\textwidth}{0.5pt} \\[0.3cm]
	{\large Anno Accademico 2026/2027}
	
\end{titlepage}
\pagenumbering{arabic}

% --- INDICE AUTOMATICO ---
\tableofcontents
\newpage

% =====================================================================
% SEZIONE 1
% =====================================================================
\section{Introduzione, Contesto Operativo e Architettura delle Infezioni}

Nell'ambito della sicurezza informatica moderna e della \textit{Threat Intelligence}, l'analisi dei campioni malevoli (\textit{Malware Analysis}) rappresenta un pilastro essenziale per estrarre Indicatori di Compromissione (IoC), comprendere i protocolli di Comando e Controllo (C2) e sviluppare regole di rilevamento efficaci. Tradizionalmente, l'indagine su un eseguibile sospetto si articola lungo due direttrici complementari: l'\textbf{analisi statica}, condotta esaminando il binario su disco mediante disassemblatori e decompilatori senza mandarlo in esecuzione, e l'\textbf{analisi dinamica}, svolta avviando il campione all'interno di un perimetro confinato --- tipicamente una macchina virtuale isolata o una \textit{sandbox} automatizzata (quale Cuckoo Sandbox o ambienti cloud analoghi) --- per osservarne in tempo reale le interazioni con il sistema operativo.

Per contrastare entrambi gli approcci e ritardare le attivit\`a di \textit{reverse engineering}, gli autori di malware contemporanei integrano sistematicamente molteplici livelli di protezione: compressione ed offuscamento dei binari (\textit{packing}), cifratura dinamica delle stringhe in memoria e, soprattutto, sofisticate routine \textbf{Anti-VM}, \textbf{Anti-Sandbox}, \textbf{Anti-Emulazione} e \textbf{Anti-Debug}. Prima di decifrare o attivare il vero carico offensivo (\textit{payload}), il codice malevolo ispeziona con estrema granularit\`a l'ambiente ospite: misura la coerenza dello scorrere del tempo durante le pause per smascherare le accelerazioni artificiali delle sandbox, accerta la presenza fisica di un operatore umano osservando il movimento del cursore e l'alternanza delle finestre grafiche, interroga l'hardware sottostante (modello e numero di core della CPU tramite WMI, dimensione della RAM e del disco rigido, traffico pregresso sulle schede di rete), verifica la connettivit\`a verso Internet e cerca tracce di strumenti diagnostici o debugger collegati al processo.

Il presente lavoro di tirocinio affronta lo studio approfondito, l'ingegnerizzazione dei meccanismi di aggiramento (\textit{bypass}) e la valutazione sperimentale dell'\textbf{emulazione binaria dinamica mirata} applicata a quattro campioni rappresentativi del panorama delle minacce reali per sistemi Microsoft Windows: \textbf{PlugX}, un \textbf{RAT in Go generico}, \textbf{Astaroth} (noto anche come \textit{Guildma}) e \textbf{Latrodectus}. Sebbene appartengano a famiglie differenti per linguaggio di sviluppo, architettura del processore e finalit\`a operative, tutti e quattro i campioni condividono un'\textbf{architettura multi-stadio} progettata per separare nettamente i controlli difensivi iniziali dall'esecuzione del modulo malevolo finale:

\begin{itemize}
	\item \textbf{Caso di Studio 1 -- PlugX (Remote Access Trojan PE32 a 32 bit in C/C++):} Articolato in due fasi distinte. Nella prima fase (\textit{Dropper e Persistenza}), l'eseguibile iniziale (\texttt{PlugXV0.exe}) sonda attivamente il sistema per escludere l'esecuzione in sandbox (verificando la coerenza del clock rispetto a pause \texttt{Sleep}, lo spostamento del mouse e la presenza di finestre in primo piano), accerta che la macchina non sia gi\`a infetta ed estrae su disco il secondo stadio configurandone l'esecuzione automatica nel Registro di Windows. Nella seconda fase (\textit{Payload RAT post-riavvio}), il binario estratto --- concepito per mascherarsi come \texttt{Media\allowbreak{}Center.\allowbreak{}exe} e protetto dal compressore eseguibile \textit{PECompact2} --- convalida il proprio percorso di lancio, inizializza la comunicazione HTTP verso il server C2 fingendosi un browser obsoleto (\textit{Internet Explorer 8}) e avvia la decompressione finale in memoria.
	
	\item \textbf{Caso di Studio 2 -- RAT in Go generico (RAT / Loader PE32+ a 64 bit compilato in Go 1.23.4):} Strutturato in tre stadi sequenziali. Nel primo stadio (\textit{Auto-rilancio nascosto}), controlla la presenza di una variabile d'ambiente privata (\texttt{\_\allowbreak{}HIDDEN\_\allowbreak{}CHILD}): non trovandola al primo lancio, genera tramite \texttt{Create\allowbreak{}ProcessW} una copia invisibile e scollegata di se stesso con la variabile valorizzata e termina immediatamente il processo padre con \texttt{os.Exit(0)} per eludere i monitoraggi focalizzati sul solo processo radice. Nel secondo stadio (\textit{Anti-VM e Ricognizione}), ispeziona la CPU tramite query WMI (\texttt{Win32\_\allowbreak{}Processor}), misura l'\textit{uptime} del sistema e la percentuale di carico del processore su un intervallo temporale, controlla RAM, volumi NTFS, traffico storico dell'interfaccia Ethernet, versione di Windows, indirizzo IP pubblico, geolocalizzazione e porte locali aperte. Nel terzo stadio (\textit{Comunicazione C2}), persiste l'identificativo univoco della vittima su disco, instaura un canale \textbf{WebSocket} (RFC 6455) verso il server remoto (\texttt{ws://\allowbreak{}54.233.190.22:8080/\allowbreak{}ws}), esfiltra l'intero profilo hardware serializzato in formato binario \textbf{Google Protocol Buffers} e resta in ascolto di comandi remoti interattivi.
	
	\item \textbf{Caso di Studio 3 -- Astaroth / Guildma (Malware multi-funzione PE32 a 32 bit in C, compresso con UPX 4.22):} Organizzato lungo cinque fasi operative. All'avvio, lo \textit{stub} del packer UPX ricostruisce in memoria RAM il codice eseguibile reale. Successivamente, il programma decifra sullo stack una stringa di controllo tramite XOR (\texttt{"Molon Labe"}) ricifrandola immediatamente in \textbf{AES-256} attraverso le Microsoft CryptoAPI (causando il crash immediato degli emulatori privi del sottosistema crittografico completo), verifica l'assenza di debugger tramite il PEB e richiede l'elevazione dei privilegi UAC (\texttt{runas}). Superati i controlli, si occulta in \texttt{App\allowbreak{}Data\textbackslash\allowbreak{}Roaming\textbackslash\allowbreak{}5kid\allowbreak{}Ro0t.\allowbreak{}exe} con attributi combinati di file nascosto e di sistema, imposta una doppia chiave di persistenza nel Registro (\texttt{HKCU} e \texttt{HKLM}), mappa l'intera sottorete locale \texttt{/24} inviando 254 pacchetti ICMP Echo Request via \textit{Raw Socket} e infine scatena un attacco distruttivo combinato (\textit{Fork-Bomb} unita alla saturazione fisica di blocchi da circa 381--400 MB di RAM a ciclo).
	
	\item \textbf{Caso di Studio 4 -- Latrodectus (Loader multi-stadio fileless PE32+ a 64 bit):} Incluso specificamente come caso di studio con esito negativo nell'emulazione dinamica user-mode pura, al fine di tracciarne con rigore scientifico i confini strutturali. Il primo stadio \`e una libreria DLL mascherata con estensione \texttt{.exe} che espone un \texttt{DllMain} intenzionalmente vuoto (per ingannare le sandbox automatiche) e nasconde il vero punto d'ingresso dietro 203 funzioni esportate identiche, protette da pesante offuscamento vettoriale SIMD, da un doppio ciclo matematico anti-emulazione da oltre 456 milioni di iterazioni e da un puntatore corrotto nell'intestazione dei moduli caricati. Il secondo stadio, cifrato nelle risorse via XOR, verifica di essere ospitato in \texttt{rundll32.exe} o \texttt{regsvr32.exe} e attua una tecnica avanzata di \textit{unhooking} e \textit{direct syscalls} leggendo \texttt{ntdll.dll} direttamente dal disco rigido per invocare il kernel scavalcando ogni intercettazione in memoria utente. Il terzo stadio (\texttt{Updater\allowbreak{}Tag.\allowbreak{}dll}) implementa infine decifratura \textit{just-in-time} delle stringhe e risoluzione dinamica delle API tramite hashing numerico.
\end{itemize}

% =====================================================================
% SEZIONE 2
% =====================================================================
\section{Analisi dei Requisiti e Specifiche del Problema}

L'obiettivo primario dello studio consiste nell'analizzare dinamicamente i campioni selezionati all'interno di un ambiente completamente isolato, deterministico e riproducibile, forzando l'esecuzione lungo tutti i rami operativi previsti dagli autori senza mai esporre un sistema operativo Windows reale n\'e consentire traffico verso infrastrutture di rete esterne. Durante la fase iniziale di triage sperimentale \`e emerso come, in assenza di contromisure specifiche, l'esecuzione dei quattro eseguibili all'interno di un emulatore binario standard si interrompa sistematicamente secondo tre distinte modalit\`a di fallimento:

\begin{enumerate}
	\item \textbf{Ibernazione stealth in cicli infiniti di attesa (PlugX e RAT in Go generico):} Un malware ben progettato evita spesso di terminare bruscamente con un codice di errore quando sospetta la presenza di una sandbox, poich\'e un arresto immediato costituirebbe di per s\'e un segnale anomalo per l'analista. In particolare, qualora \textbf{PlugX} rilevi una discrepanza tra la durata richiesta a \texttt{Sleep} e l'avanzamento restituito da \texttt{GetTickCount}, oppure riscontri coordinate del cursore immobili o l'assenza di una finestra grafica attiva, esso rientra silenziosamente all'inizio del ciclo di validazione ripetendo i controlli all'infinito e congelando l'estrazione del secondo stadio. In modo ancor pi\`u insidioso, il \textbf{RAT in Go generico} --- attraverso il proprio runtime concorrente --- gestisce le pause temporali interrogando direttamente la pagina di memoria condivisa del kernel Windows \texttt{KUSER\_\allowbreak{}SHARED\_\allowbreak{}DATA} (\texttt{0x7ffe0000}): poich\'e in un emulatore privo di timer hardware del kernel i contatori in tale pagina restano statici, la pausa di un secondo interposta tra le due letture di carico della CPU non scade mai, intrappolando il runtime Go in un \textit{deadlock} infinito.
	
	\item \textbf{Terminazione prematura del processo, crash su sottosistemi non implementati o saturazione delle risorse (RAT in Go generico e Astaroth):} Una seconda classe di ostacoli deriva dall'interazione tra tecniche evasive attive e lacune fisiologiche degli emulatori leggeri. Il \textbf{RAT in Go generico} chiude deliberatamente il processo primario subito dopo aver lanciato la copia nascosta (\texttt{\_\allowbreak{}HIDDEN\_\allowbreak{}CHILD=1}), facendo perdere immediatamente le tracce agli strumenti a singolo processo; subito dopo interroga il servizio COM/WMI di Windows (\texttt{ole32.dll}, \texttt{oleaut32.dll}), la cui assenza in emulazione provoca il fallimento fatale della routine di profilazione. Analogamente, \textbf{Astaroth} invoca all'avvio il sottosistema crittografico Microsoft (\texttt{Crypt\allowbreak{}Acquire\allowbreak{}ContextA} su \texttt{CRYPTSP.dll}), non implementato nativamente in Qiling, innescando un'eccezione di accesso invalido alla memoria prima ancora che inizi il comportamento malevolo; nella fase conclusiva, inoltre, tenta di paralizzare la macchina ospite moltiplicando esponenzialmente i processi e scrivendo fisicamente 100 milioni di interi in blocchi da 400.000.000 di byte di RAM ad ogni iterazione, causando il crash per \textit{Out-Of-Memory} (OOM) o il superamento del tempo massimo nelle sandbox.
	
	\item \textbf{Incompatibilit\`a strutturale con il paradigma di emulazione user-mode (Latrodectus e fase finale di PECompact2 in PlugX):} Alcune tecniche di protezione operano deliberatamente al di sotto del livello di astrazione delle librerie dinamiche in spazio utente. In \textbf{Latrodectus}, l'esecuzione di oltre 456 milioni di iterazioni aritmetiche e vettoriali prive di chiamate di sistema sfrutta il divario prestazionale (di diversi ordini di grandezza) tra una CPU fisica superscalare e un interprete software di istruzioni; nel secondo stadio, inoltre, la rilettura grezza di \texttt{C:\textbackslash\allowbreak{}Windows\textbackslash\allowbreak{}System32\textbackslash\allowbreak{}ntdll.\allowbreak{}dll} da disco e l'invocazione diretta dell'istruzione \texttt{syscall} da un buffer privato aggirano radicalmente qualsiasi meccanismo di \textit{API hooking} user-mode, richiedendo la presenza di un vero kernel Windows a Ring 0. Un ulteriore limite strutturale emerge quando un compressore come \textit{PECompact2} (nello Stadio 2 di PlugX) sollecita traduzioni a basso livello tra path Windows (\texttt{SysWOW64}) e chiamate POSIX del sistema ospite, portando alla luce difetti interni nel layer di astrazione dell'emulatore.
\end{enumerate}

Per rispondere rigorosamente a tali criticit\`a sono stati definiti i seguenti requisiti progettuali:
\begin{itemize}
	\item progettare per ciascun campione un \textit{harness} modulare in Python in grado di intercettare e falsificare in tempo reale i parametri temporali, di input, hardware, crittografici e di rete, mantenendo piena coerenza logica tra tutte le risposte restituite;
	\item implementare meccanismi di cattura automatica in memoria (\textit{in-memory dumping} e File System Virtuale) per salvare su disco gli stadi decifrati, i pacchetti di rete grezzi e i report telemetrici strutturati in formato JSON;
	\item affiancare all'emulazione dinamica l'analisi disassemblata e lo sviluppo di estrattori statici per recuperare e documentare integralmente anche i payload dei campioni non completabili in sola modalit\`a user-mode.
\end{itemize}

% =====================================================================
% SEZIONE 3
% =====================================================================
\section{Architettura del Sistema e Tecnologie Adottate}

\subsection{Stack Tecnologico, Isolamento Ambientale e Configurazione del \texttt{rootfs}}

La piattaforma sperimentale sviluppata nel corso del tirocinio si fonda sull'emulazione binaria dinamica mediante il framework \textbf{Qiling 1.4.6 (su Python 3.12)}, governato interamente attraverso script di controllo ed estrazione scritti in Python. Costruito al di sopra di \textit{Unicorn Engine} (derivato dal motore di traduzione dinamica della CPU di QEMU) e integrato con \textit{Capstone} per il disassemblaggio a runtime, Qiling si distingue radicalmente da un hypervisor tradizionale (come VMware, VirtualBox o KVM/QEMU full-system): anzich\'e virtualizzare l'intero hardware di un elaboratore e avviare un sistema operativo Windows completo, Qiling funge da \textit{loader} ed emulatore a livello di singolo processo \textit{user-mode}. Esso mappa direttamente nello spazio di indirizzamento virtuale simulato le sezioni dell'eseguibile Portable Executable (PE32 o PE32+), inizializza i registri del processore x86/x86\_64, costruisce in memoria le strutture fondamentali del processo Windows --- quali il \textit{Thread Environment Block} (TEB), il \textit{Process Environment Block} (PEB), la \textit{Import Address Table} (IAT) e l'heap di processo --- ed emula sequenzialmente ogni singola istruzione macchina.

Dal punto di vista della configurazione dell'ambiente di esecuzione, \`e stato predisposto un ambiente virtuale dedicato (\texttt{env}) tramite Python per isolare in modo rigoroso tutte le librerie, le dipendenze crittografiche e le installazioni del framework rispetto al sistema operativo ospite, garantendo la totale riproducibilit\`a degli esperimenti senza conflitti di versione.

Affinch\'e il loader PE di Qiling possa risolvere correttamente le tabelle di importazione degli eseguibili analizzati e mappare in memoria le librerie dinamiche di sistema, \`e indispensabile allestire un albero di directory simulato denominato \texttt{rootfs}, che riproduce la gerarchia essenziale di un volume Windows (\texttt{Windows\textbackslash\allowbreak{}System32}, \texttt{Windows\textbackslash\allowbreak{}SysWOW64}, \texttt{Users}, registro \textit{hive} di base). Per supportare con elevata fedelt\`a sia i campioni moderni a 64 bit sia gli eseguibili legacy o compressi a 32 bit, il \texttt{rootfs} di Qiling \`e stato popolato con le \textbf{librerie DLL originali di Windows 10, sia a 32 bit che a 64 bit, estratte direttamente da un'immagine ISO ufficiale} del sistema operativo; al suo interno sono state inoltre \textbf{importate le librerie DLL di Windows 7 a 32 bit}, assicurando piena copertura anche per dipendenze runtime pi\`u datate.

\subsection{Meccanismi di Intercettazione, Manipolazione della Memoria e Sottosistemi Simulati}

L'architettura degli \textit{harness} Python sviluppati nei moduli del repository (\texttt{plugx/}, \texttt{go\_rat/}, \texttt{astaroth/} e \texttt{latrodectus/}) sfrutta quattro livelli complementari di controllo messi a disposizione da Qiling:

\begin{enumerate}
	\item \textbf{API Hooking tipizzato (\texttt{@winsdkapi} e \texttt{ql.\allowbreak{}os.\allowbreak{}set\_\allowbreak{}api}):} Ogni volta che il codice del malware invoca una funzione esportata da una DLL di sistema (ad esempio \texttt{Sleep}, \texttt{Get\allowbreak{}Tick\allowbreak{}Count64}, \texttt{GetCursorPos}, \texttt{Create\allowbreak{}ProcessW}, \texttt{CryptEncrypt} o \texttt{WSASend}), Qiling intercetta il salto all'indirizzo della IAT e cede il controllo alla corrispondente funzione Python decorata con la convenzione di chiamata appropriata (\texttt{STDCALL} per le API Win32/Win64 o \texttt{CDECL} per le funzioni della libreria standard C \texttt{msvcrt.dll}). All'interno dell'hook, lo script ispeziona i parametri decodificati, legge o sovrascrive arbitrariamente buffer nella RAM simulata tramite \texttt{ql.\allowbreak{}mem.\allowbreak{}read()} e \texttt{ql.\allowbreak{}mem.\allowbreak{}write()}, aggiorna all'occorrenza il codice di errore del thread (inclusa la scrittura diretta all'offset \texttt{0x68} del TEB a 64 bit per compatibilit\`a con il runtime Go) e restituisce al chiamante un valore di ritorno sintetico pienamente conforme alle specifiche MSDN.
	
	\item \textbf{Hook Chirurgici su Indirizzi e Istruzioni (\texttt{ql.\allowbreak{}hook\_\allowbreak{}address} e \texttt{ql.\allowbreak{}hook\_\allowbreak{}code}):} Laddove il controllo o l'ostacolo non risieda in una chiamata API esterna bens\`i direttamente nel codice assembly del campione, gli script agganciano specifici indirizzi virtuali (\textit{Virtual Address}). Ci\`o consente di catturare l'esatto istante in cui uno \textit{stub} di decompressione salta all'\textit{Original Entry Point} (per effettuare il dump del binario decompattato in RAM), di forzare i registri generali (\texttt{RAX}, \texttt{EAX}) per impedire l'auto-terminazione di un processo padre, o di spostare direttamente l'\textit{Instruction Pointer} (\texttt{EIP}/\texttt{RIP}) per scavalcare cicli di centinaia di milioni di scritture progettati per mandare in stallo l'emulatore.
	
	\item \textbf{Virtualizzazione di Sottosistemi Complessi (VFS, COM/WMI e pagina condivisa \texttt{KUSER\_\allowbreak{}SHARED\_\allowbreak{}DATA}):} Per supplire all'assenza di interi sottosistemi di Windows nell'emulazione user-mode, sono stati ingegnerizzati in Python componenti software dedicati: un \textit{File System Virtuale} (VFS) basato su dizionari di buffer e cursori in memoria per gestire lo streaming C a basso livello; un emulatore di oggetti COM/OLE completo di \textit{Virtual Method Table} (VTable) binarie a 64 bit per rispondere alle interrogazioni hardware WMI (\texttt{IDispatch} e \texttt{IEnumVARIANT}); e un gestore sincronizzato dei contatori temporali ad alta risoluzione nella pagina condivisa \texttt{0x7ffe0000}.
	
	\item \textbf{Simulazione End-to-End dello Stack di Rete e dei Server C2:} Nessun pacchetto lascia mai l'ambiente isolato. Le risoluzioni DNS (\texttt{GetAddrInfoW}, \texttt{gethostbyname}), le interrogazioni HTTP verso servizi di geolocalizzazione, le scansioni su \textit{Raw Socket} ICMP e persino il server remoto dell'attaccante su protocollo \textbf{WebSocket} (RFC 6455) e \textbf{Protocol Buffers} vengono interamente emulati in memoria, permettendo di dialogare attivamente con il trojan e di decodificarne tutte le capacit\`a offensive.
\end{enumerate}

% =====================================================================
% SEZIONE 4
% =====================================================================
\section{Identificazione, Provenienza e Attribuzione dei Campioni}

Prima di esaminare nel dettaglio l'implementazione delle tecniche di evasione e le relative logiche di neutralizzazione per ciascun caso di studio, la presente sezione documenta formalmente i metadati identificativi dei quattro campioni oggetto di indagine. Al fine di garantire il pieno rigore metodologico, la tracciabilit\`a e la riproducibilit\`a sperimentale, per ogni eseguibile sono riportati l'impronta crittografica \textbf{SHA-256}, la \textbf{fonte di acquisizione} e il \textbf{criterio di attribuzione} alla rispettiva famiglia malware, elaborato incrociando i tag della comunit\`a di intelligence, i riscontri delle regole \textbf{YARA} e i verdetti dei motori antivirus su \textbf{VirusTotal}. La Tabella~\ref{tab:overview_samples} sintetizza preliminarmente il quadro comparativo del dataset analizzato.

\begin{table}[!htbp]
	\centering
	\scriptsize
	\setlength{\tabcolsep}{4.5pt}
	\renewcommand{\arraystretch}{1.35}
	\begin{tabularx}{\textwidth}{|P{1.85cm}|P{1.95cm}|P{2.20cm}|P{3.00cm}|C{0.95cm}|Y|}
		\hline
		\textbf{Campione} & \textbf{Formato PE} & \textbf{Linguaggio / Toolchain} & \textbf{Packer / Protezione} & \textbf{VT} & \textbf{Esito in Qiling 1.4.6} \\
		\hline
		\textbf{PlugX} & PE32 (x86) & C/C++ (MSVC) & Dropper custom + PECompact2 (Stage 2) & 62/65 & Bypass Stage 1, dump Stage 2 e cattura C2 HTTP; limite POSIX su \texttt{SysWOW64} \\
		\hline
		\textbf{RAT in Go generico} & PE32+ (x64) & Go 1.23.4 & Offuscamento runtime Go + spawn nascosto & 21/70 & Emulazione completa Stage 1--3 e ciclo interattivo C2 WebSocket/\allowbreak{}Protobuf \\
		\hline
		\textbf{Astaroth} & PE32 (x86) & C (MinGW GCC 13) & UPX 4.22 + XOR \texttt{0x5A} / AES-256 & 42/56 & Unpacking automatico in RAM e completamento fino a terminazione \\
		\hline
		\textbf{Latrodectus} & PE32+ DLL (\texttt{.exe}) & C/C++ (x86\_64) & Loop matematico SIMD + XOR + Direct Syscall & 44/70 & Caso negativo user-mode; estrazione statica completa Stadi 2 e 3 \\
		\hline
	\end{tabularx}
	\caption{Quadro sinottico dei quattro campioni analizzati, delle relative caratteristiche binarie e dell'esito sperimentale ottenuto.}
	\label{tab:overview_samples}
\end{table}

\subsection{Caso di Studio 1: PlugX}

\begin{itemize}
	\item \textbf{SHA-256:} {\small\texttt{7362590eec1693d6ead5f274aa546f276bcc6939962f195eb997fcd57f02be02}}
	\item \textbf{Fonte:} MalwareBazaar
	\item \textbf{Tag:} \texttt{Plugx}
	\item \textbf{Regole YARA rilevate:}
	\begin{itemize}
		\item \textit{Offuscamento:} \texttt{malware\_\allowbreak{}sakula\_\allowbreak{}xorloop} (autore: David Cannings).
		\item \textit{Famiglia / Payload:} \texttt{RAT\_Sakula} (autori: Airbus Defence and Space Cybersecurity CSIRT -- Yoann Francou / NCC Group David Cannings); \texttt{Scan\allowbreak{}Box\_\allowbreak{}Malware\_\allowbreak{}Generic} (autore: Florian Roth).
		\item \textit{Packer:} \texttt{suspicious\_\allowbreak{}packer\_\allowbreak{}section} (autore: \texttt{@j0sm1}).
	\end{itemize}
	\item \textbf{Rilevamenti Antivirus:} 62/65 su VirusTotal. Tra i principali security vendor: Microsoft (\texttt{Backdoor:Win32/\allowbreak{}Plugx!pz}) e QuickHeal (\texttt{Backdoor.\allowbreak{}PlugX.\allowbreak{}S13515199}).
	\item \textbf{Criterio di Attribuzione:} Famiglia PlugX, confermata univocamente dal tag su MalwareBazaar e dall'ampio consenso dei rilevamenti antivirus (62/65); i riscontri YARA riconducibili a Sakula e ScanBox evidenziano il noto riutilizzo di routine di offuscamento XOR e di loader multi-stadio condivisi all'interno dell'ecosistema di minacce APT affini.
\end{itemize}

\subsection{Caso di Studio 2: RAT in Go generico}

\begin{itemize}
	\item \textbf{SHA-256:} {\small\texttt{818a48e399389a23fc73994dc4331314374a7097fe831b4ee617b9315487d804}}
	\item \textbf{Fonte:} MalwareBazaar
	\item \textbf{Tag:} \texttt{54-\allowbreak{}233-\allowbreak{}190-\allowbreak{}22}, \texttt{exe}
	\item \textbf{Regole YARA rilevate:}
	\begin{itemize}
		\item \textit{Anti-Debugging:} \texttt{Debugger\allowbreak{}Check\_\_\allowbreak{}Query\allowbreak{}Info}, \texttt{Debugger\allowbreak{}Exception\_\_\allowbreak{}Console\allowbreak{}Ctrl}, \texttt{Debugger\allowbreak{}Exception\_\_\allowbreak{}Set\allowbreak{}Console\allowbreak{}Ctrl}.
		\item \textit{Compilazione / Linguaggio:} \texttt{Detect\allowbreak{}Go\allowbreak{}Method\allowbreak{}Signatures} (autore: Wyatt Tauber), \texttt{GoBinTest}, \texttt{golang}, \texttt{Golangmalware}.
		\item \textit{Famiglia:} \texttt{Hive\allowbreak{}Ransomware} (autore: Dhanunjaya), \texttt{RANSOMWARE} (autore: ToroGuitar).
		\item \textit{Offuscamento / Struttura PE:} \texttt{Sus\_\allowbreak{}Obf\_\allowbreak{}Enc\_\allowbreak{}Spoof\_\allowbreak{}Hide\_\allowbreak{}PE} (autore: XiAnzheng), \texttt{Thread\allowbreak{}Control\_\_\allowbreak{}Context}.
	\end{itemize}
	\item \textbf{Rilevamenti Antivirus:} 21/70 su VirusTotal. Tra i principali security vendor: Microsoft (\texttt{Trojan:Win32/\allowbreak{}Yomal!rfn}) e Kaspersky (\texttt{Hack\allowbreak{}Tool.\allowbreak{}Win32.\allowbreak{}GO\allowbreak{}Simple\allowbreak{}Tunnel.\allowbreak{}as}).
	\item \textbf{Criterio di Attribuzione:} Classificato come \textbf{RAT in Go generico} poich\'e non vi \`e alcuna conferma n\'e dai tag su MalwareBazaar (limitati all'indirizzo IP del C2 \texttt{54.233.190.22} e al formato eseguibile \texttt{exe}) n\'e dai verdetti dei security vendor circa l'appartenenza a una famiglia nota. Le firme YARA intercettano prevalentemente le strutture standard del compilatore e del runtime Go (circostanza che genera falsi positivi verso altre famiglie scritte nel medesimo linguaggio, come \texttt{Hive\allowbreak{}Ransomware}), mentre i motori antivirus riportano rilevazioni generiche (\texttt{Trojan:Win32/\allowbreak{}Yomal!rfn}) o riconducibili a utilit\`a di tunneling in Go (\texttt{GO\allowbreak{}Simple\allowbreak{}Tunnel}), in piena coerenza con la reale natura di \textit{Remote Access Trojan} custom emersa dall'analisi dinamica.
\end{itemize}

\subsection{Caso di Studio 3: Astaroth}

\begin{itemize}
	\item \textbf{SHA-256:} {\small\texttt{9b0d5e40ea39bcdb7f21c195750b010d7ebe343eefd691b27e572e6bbd740c33}}
	\item \textbf{Fonte:} MalwareBazaar
	\item \textbf{Tag:} \texttt{Astaroth}, \texttt{Evader}, \texttt{exe}, \texttt{forkboom}, \texttt{guildma}, \texttt{Locker}, \texttt{UPX}
	\item \textbf{Regole YARA rilevate:}
	\begin{itemize}
		\item \textit{Packer / Offuscamento:} \texttt{MAL\_\allowbreak{}packer\_\allowbreak{}lb\_\allowbreak{}was\_\allowbreak{}detected} (autore: \texttt{0x0d4y}), \texttt{Sus\_\allowbreak{}Obf\_\allowbreak{}Enc\_\allowbreak{}Spoof\_\allowbreak{}Hide\_\allowbreak{}PE} (autore: XiAnzheng), \texttt{UPX290LZMA\allowbreak{}Markus\allowbreak{}Oberhumer\allowbreak{}Laszlo\allowbreak{}Molnar\allowbreak{}John\allowbreak{}Reiser} (autore: \texttt{malware-lu}).
		\item \textit{Struttura PE:} \texttt{pe\_\allowbreak{}detect\_\allowbreak{}tls\_\allowbreak{}callbacks}.
	\end{itemize}
	\item \textbf{Rilevamenti Antivirus:} 42/56 su VirusTotal. Tra i principali security vendor: Microsoft (\texttt{Trojan:Win32/\allowbreak{}Egairtigado!rfn}), Malwarebytes (\texttt{Trojan.Crypt}) e McAfee Scanner (\texttt{Ti!9B0D5E40EA39}).
	\item \textbf{Criterio di Attribuzione:} Famiglia Astaroth / Guildma, confermata dai tag descrittivi su MalwareBazaar (incluso \texttt{forkboom} e \texttt{UPX}, pienamente riscontrati nell'analisi dinamica) nonostante i rilevamenti statici YARA e antivirus si fermino prevalentemente allo strato esterno di compressione UPX.
\end{itemize}

\subsection{Caso di Studio 4: Latrodectus}

\begin{itemize}
	\item \textbf{SHA-256:} {\small\texttt{fc58ae120d35b751525f6c312f9aeff75e9151f02fa113815bd01acc49f63596}}
	\item \textbf{Fonte:} MalwareBazaar
	\item \textbf{Tag:} \texttt{exe}, \texttt{Latrodectus}
	\item \textbf{Regole YARA rilevate:}
	\begin{itemize}
		\item \textit{Anti-Debugging:} \texttt{Debugger\allowbreak{}Check\_\_\allowbreak{}API}.
		\item \textit{Offuscamento:} \texttt{Detect\allowbreak{}Encrypted\allowbreak{}Variants} (autore: Zinyth).
		\item \textit{Compilazione / Linguaggio:} \texttt{golang}, \texttt{golang\_\allowbreak{}bin\_\allowbreak{}JCorn\_\allowbreak{}CSC846} (autore: Justin Cornwell), \texttt{identity\_\allowbreak{}golang} (autore: Eric Yocam).
	\end{itemize}
	\item \textbf{Rilevamenti Antivirus:} 44/70 su VirusTotal. Tra i principali security vendor: Kaspersky (\texttt{Backdoor.\allowbreak{}Win64.\allowbreak{}Latrodectus.\allowbreak{}ar}), Lionic (\texttt{Trojan.\allowbreak{}Win32.\allowbreak{}Latrodectus.\allowbreak{}m!c}) e AliCloud (\texttt{Backdoor:Win/\allowbreak{}Latrodectus.\allowbreak{}ax}).
	\item \textbf{Criterio di Attribuzione:} Famiglia Latrodectus, confermata in modo convergente sia dai tag su MalwareBazaar sia dalle firme specifiche dei motori antivirus (Kaspersky, Lionic, AliCloud).
\end{itemize}

% =====================================================================
% SEZIONE 5
% =====================================================================
\section{Caso di Studio 1: Analisi Multi-Stadio e Bypass Anti-VM nel Malware PlugX}

Il primo caso di studio riguarda \textbf{PlugX}, uno storico e longevo \textit{Remote Access Trojan} impiegato in numerose campagne di spionaggio informatico. L'ispezione del campione iniziale (\texttt{PlugXV0.exe}, eseguibile PE32 a 32 bit scritto in C/C++) mediante l'harness sviluppato in \texttt{plugx\_\allowbreak{}multistage\_\allowbreak{}unpacker\_\allowbreak{}analyzer.\allowbreak{}py} ha messo in luce un'architettura scissa in due stadi distinti: un \textbf{Dropper di Fase 1}, responsabile della validazione ambientale Anti-Sandbox, della verifica di non-reinfezione e dell'estrazione su disco del componente successivo, e un \textbf{Payload RAT di Fase 2} (\texttt{Media\allowbreak{}Center.\allowbreak{}exe}), compresso tramite \textit{PECompact2} e progettato per entrare in azione esclusivamente al riavvio del sistema operativo dalla directory di persistenza.

\subsection{Prima Fase: Evasione Ambientale ed Estrazione del Secondo Stadio (\texttt{PlugXV0.exe})}

All'avvio del primo stadio, il malware non decifra immediatamente il payload incorporato, ma entra in una gabbia di validazione ambientale che interroga ciclicamente il sistema operativo. Qualora anche uno solo dei controlli fallisca, PlugX evita di terminare il processo --- comportamento che insospettirebbe l'analista --- e ripete indefinitamente il ciclo di verifica, ibernando di fatto l'esecuzione. Per mantenere uno stato simulato internamente coerente attraverso tutte le invocazioni, nello script Python \`e stata progettata la classe \texttt{AntiEvasion}, incaricata di coordinare orologio interno, coordinate del puntatore, handle delle finestre e identificativi hardware.

\subsubsection{Alterazione Temporale (\texttt{Sleep} e \texttt{GetTickCount})}

\textbf{Funzionamento della Tecnica Anti-Sandbox:} Gli ambienti automatizzati di analisi dinamica sono vincolati da rigide finestre temporali (spesso pochi minuti per campione) e tendono pertanto a intercettare le funzioni di attesa come \texttt{Sleep} ignorandole (\textit{sleep skipping}) o contraendone drasticamente la durata. PlugX smaschera questa manipolazione invocando l'API \texttt{GetTickCount} (che restituisce i millisecondi trascorsi dall'avvio di Windows) immediatamente prima e immediatamente dopo una pausa di 5.000 millisecondi (\texttt{Sleep(5000)}). Se la differenza $\Delta t = t_2 - t_1$ calcolata dal malware risulta inferiore alla soglia dei cinque secondi attesi, la routine deduce che il tempo \`e stato alterato artificialmente da una sandbox e ricomincia da capo il ciclo difensivo.

\textbf{Logica di Risoluzione e Implementazione:} Per azzerare i tempi morti dell'analisi senza far fallire il controllo aritmetico del malware, lo script intercetta entrambe le API disaccoppiando il tempo reale dal tempo percepito dal binario. Quando il processo chiama \texttt{Sleep(dw\allowbreak{}Milliseconds)}, l'hook ritorna istantaneamente \texttt{0} (annullando qualsiasi attesa reale), ma accumula i millisecondi richiesti nell'attributo \texttt{tempo\_\allowbreak{}in\_\allowbreak{}sospeso}. Alla successiva chiamata a \texttt{GetTickCount}, l'emulatore somma al proprio contatore interno sia un lieve avanzamento fisiologico della CPU (\texttt{+50 ms}), sia l'intero credito temporale parcheggiato in precedenza, azzerando poi l'accumulatore (Listato~\ref{lst:plugx_timing}).

\begin{lstlisting}[style=pythoncode, caption={Gestione coerente dell'aritmetica temporale per il bypass dello Sleep-Skipping in PlugX.}, label={lst:plugx_timing}]
@winsdkapi(cc=STDCALL, params={"dwMilliseconds": DWORD})
def hook_Sleep(ql, address, params):
    tempo_richiesto = params["dwMilliseconds"]
    anti_evasion.tempo_in_sospeso += tempo_richiesto
    print(f"[DEBUG TEMPO] Salvati {tempo_richiesto} ms (Sleep bypassato).")
    return 0

@winsdkapi(cc=STDCALL, params={})
def hook_GetTickCount(ql, address, params):
    # Avanza di 50 ms fisiologici + eventuali millisecondi accumulati da Sleep()
    anti_evasion.tick += 50 + anti_evasion.tempo_in_sospeso
    anti_evasion.tempo_in_sospeso = 0
    return anti_evasion.tick
\end{lstlisting}

In questo modo una pausa nominale di cinque secondi viene superata in meno di un millisecondo effettivo sulla macchina di analisi, ma la sottrazione eseguita dal malware rileva esattamente \texttt{5050 ms} trascorsi, convalidando la genuinit\`a dell'orologio di sistema.

\subsubsection{Controllo del Movimento del Cursore (\texttt{GetCursorPos})}

\textbf{Funzionamento della Tecnica Anti-VM:} Nelle macchine virtuali presidiate da sistemi automatici non vi \`e alcun operatore fisico che muova il mouse. PlugX sfrutta questa evidenza invocando due volte l'API \texttt{GetCursorPos} (separando le due letture con un intervallo temporale) per popolare in memoria una struttura \texttt{POINT} contenente le coordinate intere con segno \texttt{(x, y)} a 32 bit ciascuna. Il codice assembly confronta quindi le coordinate orizzontali registrate: se la sottrazione restituisce zero (cursore immobile), il campione rimane bloccato nel loop conservativo.

\textbf{Logica di Risoluzione e Implementazione:} Agganciando \texttt{GetCursorPos}, l'harness Python genera ad ogni invocazione una micro-oscillazione pseudocasuale (\textit{jitter}) compresa nell'intervallo $[-15, +15]$ pixel su entrambi gli assi a partire dalla posizione iniziale \texttt{(500, 500)}, impacchettando i due interi a 32 bit in formato \textit{little-endian} direttamente all'indirizzo puntato da \texttt{lpPoint} tramite \texttt{ql.mem.write} (Listato~\ref{lst:plugx_hid}).

\subsubsection{Controllo della Finestra Attiva (\texttt{GetForegroundWindow}) e Difetto Logico Interno}

\textbf{Funzionamento della Tecnica e Scoperta della Vulnerabilit\`a:} Per verificare l'esistenza di una sessione desktop interattiva, PlugX interroga \texttt{Get\allowbreak{}Foreground\allowbreak{}Window}, che restituisce l'identificativo (\texttt{HWND}) della finestra attualmente in primo piano. Negli ambienti \textit{headless} o mal configurati la chiamata restituisce \texttt{0}, bloccando immediatamente il malware. L'analisi statica condotta in parallelo tramite disassemblaggio (e confermata dal tracciatore assembly dinamico integrato nello script tra gli indirizzi \texttt{0x004012E0} e \texttt{0x004012F2}) ha tuttavia portato alla luce una significativa negligenza implementativa commessa dagli sviluppatori del trojan: quando \texttt{Get\allowbreak{}Foreground\allowbreak{}Window} restituisce un handle non nullo, il binario lo memorizza correttamente per confrontarlo in seguito e verificare che l'utente cambi finestra nel tempo; tuttavia, a causa di una variabile flag interna inizializzata staticamente a zero e mai aggiornata nel ramo principale, il flusso incontra un salto incondizionato che scavalca per intero la seconda verifica comparativa.

\textbf{Logica di Risoluzione e Implementazione:} Facendo leva su tale difetto logico del campione, \`e risultato sufficiente intercettare \texttt{Get\allowbreak{}Foreground\allowbreak{}Window} e restituire un handle sintetico mai nullo e progressivamente crescente (\texttt{0x10004}, \texttt{0x10008}, \dots), completando al contempo la profilazione desktop con risoluzioni video realistiche (\texttt{1920x1080} in \texttt{Get\allowbreak{}System\allowbreak{}Metrics}), nome host credibile (\texttt{DESKTOP-\allowbreak{}OFFICE} in \texttt{Get\allowbreak{}Computer\allowbreak{}NameA}), seriale di volume valido (\texttt{0x9A2B3C4D} in \texttt{Get\allowbreak{}Volume\allowbreak{}InformationA}) e privilegi amministrativi attivi (\texttt{Is\allowbreak{}User\allowbreak{}An\allowbreak{}Admin -> 1}), come mostrato nel Listato~\ref{lst:plugx_hid}.

\begin{lstlisting}[style=pythoncode, caption={Simulazione dell'interazione umana (HID/GUI) su GetCursorPos e GetForegroundWindow.}, label={lst:plugx_hid}]
@winsdkapi(cc=STDCALL, params={"lpPoint": POINTER})
def hook_GetCursorPos(ql, address, params):
    anti_evasion.cursor_x += random.randint(-15, 15)
    anti_evasion.cursor_y += random.randint(-15, 15)
    ql.mem.write(
        params["lpPoint"],
        anti_evasion.cursor_x.to_bytes(4, "little", signed=True)
        + anti_evasion.cursor_y.to_bytes(4, "little", signed=True),
    )
    return 1

@winsdkapi(cc=STDCALL, params={})
def hook_GetForegroundWindow(ql, address, params):
    anti_evasion.window_handle += 4
    return anti_evasion.window_handle
\end{lstlisting}

\subsubsection{Controllo di Pre-Infezione ed Estrazione tramite File System Virtuale (VFS)}

\textbf{Funzionamento della Tecnica:} Superata la barriera dei controlli ambientali, il dropper esegue un test di pre-infezione aprendo in modalit\`a sola lettura (\texttt{fopen} con modo \texttt{"rb"}) il percorso destinato al secondo stadio: \texttt{C:\textbackslash\allowbreak{}Users\textbackslash\allowbreak{}Admin\textbackslash\allowbreak{}App\allowbreak{}Data\textbackslash\allowbreak{}Local\textbackslash\allowbreak{}Temp\textbackslash\allowbreak{}Micro\allowbreak{}Media\textbackslash\allowbreak{}Media\allowbreak{}Center.\allowbreak{}exe} (costruito dinamicamente espandendo \texttt{\%TEMP\%} tramite \texttt{Expand\allowbreak{}Environment\allowbreak{}StringsA}). Se l'apertura ha successo (puntatore \texttt{FILE*} diverso da zero), il malware conclude che la macchina \`e gi\`a compromessa e termina immediatamente per evitare conflitti o esecuzioni duplicate. In caso negativo, riapre il proprio eseguibile su disco (\texttt{PlugXV0.exe}), si posiziona tramite \texttt{fseek} all'offset dell'overlay cifrato, legge il blocco dati con \texttt{fread}, lo decifra in memoria RAM e lo riversa sul nuovo file tramite \texttt{fwrite}.

\textbf{Logica di Risoluzione e Implementazione del Custom VFS:} Durante i test in Qiling \`e emerso che le primitive di I/O bufferizzato della libreria standard C (\texttt{msvcrt.dll}: \texttt{fopen}, \texttt{\_fsopen}, \texttt{fseek}, \texttt{fread}, \texttt{fwrite}, \texttt{ftell}, \texttt{fclose}, \texttt{\_lock}, \texttt{\_unlock}) non erano pienamente supportate dal layer nativo dell'emulatore, generando letture su memoria non mappata. Si \`e quindi progettato in Python un \textbf{File System Virtuale (VFS)} completo in RAM basato sui dizionari \texttt{file\_\allowbreak{}buffers} e \texttt{file\_\allowbreak{}cursors} indicizzati da puntatori allocati sull'heap simulato (Listato~\ref{lst:plugx_vfs}):
\begin{itemize}
	\item quando \texttt{fopen} richiede in lettura \texttt{Media\allowbreak{}Center.\allowbreak{}exe} durante lo Stadio 1, il VFS restituisce \texttt{0} (\texttt{NULL}), convincendo il dropper che il sistema sia pulito;
	\item quando \texttt{fopen} apre \texttt{PlugXV0.exe}, il VFS carica in memoria i byte reali del campione e gestisce con precisione aritmetica tutti gli spostamenti di \texttt{fseek} (\texttt{SEEK\_SET}, \texttt{SEEK\_CUR}, \texttt{SEEK\_END}) e i trasferimenti in RAM di \texttt{fread};
	\item non appena il malware invoca \texttt{fwrite} per salvare il secondo stadio appena decifrato, l'hook intercetta il buffer sorgente nella memoria di Qiling e lo esporta direttamente sul disco dell'analista come \texttt{dropped\_\allowbreak{}payload\_\allowbreak{}0x5014f15.\allowbreak{}bin}, catturando integralmente i \textbf{74.358 byte} del payload RAT.
\end{itemize}

\begin{lstlisting}[style=pythoncode, caption={Estratto del File System Virtuale (VFS) per il bypass di pre-infezione e il dump di MediaCenter.exe.}, label={lst:plugx_vfs}]
@winsdkapi(cc=CDECL, params={"filename": STRING, "mode": STRING})
def hook_fopen(ql, address, params):
    filename = (params["filename"] or "")
    mode = (params["mode"] or "")
    ptr = ql.os.heap.alloc(0x1000)

    if STADIO_ANALISI == 1:
        # Bypass controllo di pre-infezione: simula l'assenza di MediaCenter.exe
        if "r" in mode and "mediacenter.exe" in filename.lower():
            return 0
        # Carica nel VFS i byte del dropper originale per consentire fseek/fread
        if "plugxv0.exe" in filename.lower():
            with open(_active_sample_path, "rb") as f:
                file_buffers[ptr] = f.read()
            file_cursors[ptr] = 0
            return ptr

    file_buffers[ptr] = b""
    file_cursors[ptr] = 0
    return ptr

@winsdkapi(cc=CDECL, params={"buffer": POINTER, "size": SIZE_T, "count": SIZE_T, "stream": POINTER})
def hook_fwrite(ql, address, params):
    stream = params["stream"]
    total_bytes = params["size"] * params["count"]
    if total_bytes > 0:
        data = ql.mem.read(params["buffer"], total_bytes)
        dump_path = os.path.join(_output_dir, f"dropped_payload_{hex(stream)}.bin")
        write_mode = "ab" if dump_path in _initialized_dumps else "wb"
        _initialized_dumps.add(dump_path)
        with open(dump_path, write_mode) as f:
            f.write(data)
        print(f"[+] DUMP COMPLETATO: {total_bytes} bytes salvati in '{dump_path}'")
    return params["count"]
\end{lstlisting}

\subsection{Seconda Fase: Analisi Post-Riavvio, Rete C2 e Limiti su \texttt{PECompact2}}

Estratto il secondo stadio (\texttt{dropped\_\allowbreak{}payload\_\allowbreak{}0x5014f15.\allowbreak{}bin}, eseguibile PE32 compresso con \textit{PECompact2}), l'harness \`e stato avviato in modalit\`a post-riavvio (\texttt{--stage 2}) per studiarne in isolamento il comportamento operativo.

\subsubsection{Falsificazione del Percorso di Esecuzione (\texttt{GetModuleFileNameA})}

Avviato fuori dal contesto di infezione originario, il secondo stadio interroga immediatamente \texttt{Get\allowbreak{}Module\allowbreak{}File\allowbreak{}NameA} per verificare la propria ubicazione sul file system: se il percorso corrente non coincide con la directory di persistenza configurata nella chiave di avvio automatico \texttt{Run} (\texttt{C:\textbackslash\allowbreak{}Users\textbackslash\allowbreak{}Admin\textbackslash\allowbreak{}App\allowbreak{}Data\textbackslash\allowbreak{}Local\textbackslash\allowbreak{}Temp\textbackslash\allowbreak{}Micro\allowbreak{}Media\textbackslash\allowbreak{}Media\allowbreak{}Center.\allowbreak{}exe}), l'esecuzione termina silenziosamente per impedire l'analisi del file rinominato su una postazione di laboratorio. Intercettando \texttt{Get\allowbreak{}Module\allowbreak{}File\allowbreak{}NameA} e iniettando direttamente nel buffer del processo la stringa attesa (Listato~\ref{lst:plugx_stage2}), il controllo di posizione \`e stato aggirato con successo.

\begin{lstlisting}[style=pythoncode, caption={Path spoofing post-riavvio e intercettazione del beaconing HTTP nello Stadio 2 di PlugX.}, label={lst:plugx_stage2}]
@winsdkapi(cc=STDCALL, params={"hModule": HANDLE, "lpFilename": POINTER, "nSize": DWORD})
def hook_GetModuleFileNameA(ql, address, params):
    fake_path = b"C:\\Users\\Admin\\AppData\\Local\\Temp\\MicroMedia\\MediaCenter.exe\x00"
    ql.mem.write(params["lpFilename"], fake_path)
    return len(fake_path) - 1

@winsdkapi(cc=STDCALL, params={
    "lpszAgent": STRING, "dwAccessType": DWORD,
    "lpszProxy": STRING, "lpszProxyBypass": STRING, "dwFlags": DWORD
})
def hook_InternetOpenA(ql, address, params):
    print(f"[BEHAVIOR/NETWORK] InternetOpenA (User-Agent: {params['lpszAgent']})")
    return 0x12345678
\end{lstlisting}

\subsubsection[Intercettazione del Traffico C2 e Limite su SysWOW64]{Intercettazione del Traffico C2 e Blocco Interno dell'Emulatore su \texttt{SysWOW64}}

Convalidato il percorso di lancio, il campione ha mappato le librerie di rete Windows (\texttt{ws2\_32.dll} e \texttt{wininet.dll}), invocando \texttt{WSAStartup} e preparando la sessione HTTP verso il server di Comando e Controllo tramite \texttt{Internet\allowbreak{}OpenA}. L'output catturato dall'harness ha svelato come il RAT camuffi il proprio traffico di \textit{beaconing} adottando un'intestazione \texttt{User-Agent} statica che simula \textit{Internet Explorer 8} su Windows XP (\texttt{Windows NT 5.1}).

Subito dopo l'inizializzazione di rete, lo stub del compressore \textit{PECompact2} ha avviato la routine di decompressione dell'ultimo strato destinato all'iniezione in memoria (per la quale l'harness aveva predisposto gli hook su \texttt{Create\allowbreak{}ProcessA} con flag \texttt{CREATE\_\allowbreak{}SUSPENDED = 0x4}, \texttt{Virtual\allowbreak{}Alloc\allowbreak{}Ex}, \texttt{Write\allowbreak{}Process\allowbreak{}Memory} e \texttt{ResumeThread}). Prima di procedere, tuttavia, il packer ha verificato l'integrit\`a dell'ambiente a 32 bit su sistemi a 64 bit richiedendo il percorso \texttt{SysWOW64} tramite \texttt{Get\allowbreak{}System\allowbreak{}Wow64\allowbreak{}DirectoryA} e ispezionando la cartella con \texttt{Find\allowbreak{}First\allowbreak{}FileA}. Nonostante l'hook Python avesse restituito correttamente \texttt{C:\textbackslash\allowbreak{}Windows\textbackslash\allowbreak{}SysWOW64} e la directory fosse presente nel \texttt{rootfs}, l'emulazione si \`e arrestata a causa di un limite strutturale interno di Qiling (Listato~\ref{lst:plugx_out}): nel tradurre l'ispezione della directory Windows sulle primitive POSIX del sistema ospite Linux, l'implementazione interna di \texttt{Find\allowbreak{}First\allowbreak{}FileA} di Qiling tenta erroneamente di aprire il percorso della cartella come un file regolare, sollevando l'eccezione Python irreversibile \texttt{IsADirectory\allowbreak{}Error: [Errno 21]}.


\noindent\begin{minipage}{\linewidth}
\begin{lstlisting}[style=consoleout, caption={Traccia di esecuzione dello Stadio 2 di PlugX: path spoofing, cattura User-Agent C2 ed eccezione su SysWOW64.}, label={lst:plugx_out}]
[BEHAVIOR/EVASION] GetModuleFileNameA falsato! (STADIO 2) Il malware crede di essere in: C:\Users\Admin\AppData\Local\Temp\MicroMedia\MediaCenter.exe
[BEHAVIOR/NETWORK] WSAStartup: Il malware sta inizializzando le librerie di rete!
[BEHAVIOR/NETWORK] InternetOpenA: Preparazione connessione HTTP (User-Agent: Mozilla/4.0+(compatible;+MSIE+8.0;+Windows+NT+5.1;+SV1))
[DEBUG/FIX] GetSystemWow64DirectoryA richiesto dal packer -> Fornito: C:\Windows\SysWOW64
IsADirectoryError: [Errno 21] Is a directory: '/home/remnux/.../Windows/SysWOW64'
[!] Eccezione durante l'esecuzione: Windows API Implementation Error
\end{lstlisting}
\end{minipage}


\subsection{Sintesi Comportamentale di PlugX}

L'indagine combinata sui due stadi ha permesso di ricostruire l'intera catena operativa del campione PlugX:
\begin{enumerate}
	\item \textbf{Validazione Ambientale e Drop:} Accerta la coerenza temporale su pause di 5 secondi, il movimento fisico del mouse e la presenza di finestre attive, verifica l'assenza di precedenti installazioni e riversa nella cartella temporanea l'eseguibile \texttt{Media\allowbreak{}Center.\allowbreak{}exe} (74.358 byte).
	\item \textbf{Persistenza e Anti-Forensics:} Configura l'esecuzione automatica nel Registro di Windows e predispone la rimozione ritardata del dropper iniziale per cancellare le tracce dell'infezione primaria.
	\item \textbf{Attivazione Post-Riavvio e Beaconing HTTP:} Al riavvio successivo, convalida l'esecuzione da \texttt{Micro\allowbreak{}Media\textbackslash\allowbreak{}Media\allowbreak{}Center.\allowbreak{}exe}, inizializza il canale HTTP verso il server C2 camuffandosi da Internet Explorer 8 e tenta lo spacchettamento finale protetto da \textit{PECompact2}.
\end{enumerate}

% =====================================================================
% SEZIONE 6
% =====================================================================
\section{Caso di Studio 2: Analisi Multi-Stadio, Runtime Go e Simulazione C2 WebSocket/Protobuf nel RAT in Go Generico}

Il secondo caso di studio affronta un campione notevolmente pi\`u complesso dal punto di vista architetturale: un \textbf{RAT in Go generico} (\texttt{go\_rat.exe}), eseguibile PE32+ a 64 bit compilato in linguaggio \textbf{Go 1.23.4} e analizzato mediante l'harness completo \texttt{gorat\_\allowbreak{}multistage\_\allowbreak{}c2\_\allowbreak{}analyzer.\allowbreak{}py}. A differenza dei tradizionali binari scritti in C/C++, un eseguibile Go incorpora staticamente un voluminoso \textit{runtime} concorrente dotato di scheduler cooperativo/preemptive di \textit{goroutine}, \textit{Garbage Collector} parallelo, gestione dedicata dei segnali asincroni, allocatore virtuale proprio e networking basato su \textit{I/O Completion Ports} (IOCP). L'analisi statica e dinamica ha rivelato un ciclo di vita strutturato in tre fasi: auto-rilancio invisibile tramite variabile d'ambiente privata, ricognizione intensiva dell'hardware e della rete in funzione Anti-VM/Anti-Sandbox e, infine, comunicazione bidirezionale persistente su protocollo \textbf{WebSocket} (RFC 6455) con payload binari serializzati in \textbf{Google Protocol Buffers} verso il server C2 \texttt{54.233.190.22:8080}.

\subsection{Adattamento Chirurgico del Runtime Go 1.23 a 64 bit in Emulazione Single-Thread}

Prima ancora di raggiungere il codice applicativo del malware (\texttt{main.main}), l'avvio di un binario Go a 64 bit dentro un emulatore \textit{user-mode} a singolo flusso di esecuzione come Qiling si scontra con le assunzioni strutturali del runtime Go. Per stabilizzare l'ambiente senza alterare la semantica del programma sono stati progettati quattro interventi di basso livello (Listato~\ref{lst:gorat_goruntime}):

\begin{enumerate}
	\item \textbf{Bypass del Garbage Collector concorrente (\texttt{runtime.\allowbreak{}gcenable} a \texttt{0x417560}):} L'attivazione delle goroutine di background del GC in assenza di uno scheduler multi-thread del kernel generava stalli durante le transizioni di stato; intercettando \texttt{0x417560} e simulando un ritorno immediato (\texttt{pop RIP}), la memoria viene preservata linearmente per l'intera durata dell'analisi.
	\item \textbf{Neutralizzazione di \texttt{os/\allowbreak{}signal.\allowbreak{}loop} (\texttt{0x53b240}):} Il thread secondario preposto alla ricezione dei segnali di sistema causava in single-thread il \textit{panic} fatale \texttt{signal\_\allowbreak{}recv: inconsistent state}; la scrittura in memoria di un'istruzione \texttt{RET} (\texttt{0xC3}) all'indirizzo \texttt{0x53b240} ha eliminato alla radice l'anomalia.
	\item \textbf{Sincronizzazione del codice di errore nel TEB a 64 bit (\texttt{GS:[0x68]}):} A differenza dei programmi C che interrogano l'API \texttt{Get\allowbreak{}Last\allowbreak{}Error()}, molte primitive della libreria standard Go leggono direttamente il campo \texttt{Last\allowbreak{}Error\allowbreak{}Value} all'offset \texttt{0x68} del \textit{Thread Environment Block}. La funzione ausiliaria \texttt{set\_\allowbreak{}last\_\allowbreak{}error(ql, code)} aggiorna simultaneamente sia l'attributo interno di Qiling sia la memoria fisica in \texttt{teb\_\allowbreak{}addr + 0x68}, garantendo che codici cruciali come \texttt{WSA\_\allowbreak{}IO\_\allowbreak{}PENDING} (\texttt{997}) o \texttt{ERROR\_\allowbreak{}FILE\_\allowbreak{}NOT\_\allowbreak{}FOUND} (\texttt{2}) vengano recepiti fedelmente dal runtime Go.
	\item \textbf{Incapsulamento degli Handle I/O (\texttt{Ql\allowbreak{}Virtual\allowbreak{}File\allowbreak{}Handle}):} Per evitare eccezioni interne nell'\texttt{Handle\allowbreak{}Manager} di Qiling durante la virtualizzazione di dispositivi speciali Windows (\texttt{NUL}, \texttt{CONOUT\$}) e delle operazioni di posizionamento a 64 bit (\texttt{Set\allowbreak{}File\allowbreak{}Pointer\allowbreak{}Ex}), gli stream Python sono stati avvolti in un wrapper dotato degli attributi \texttt{.id}, \texttt{.permissions} e \texttt{.inherit} attesi da \texttt{CloseHandle}.
\end{enumerate}

\begin{lstlisting}[style=pythoncode, caption={Adattamento del runtime Go 1.23 a 64 bit e sincronizzazione diretta del LastError nel TEB.}, label={lst:gorat_goruntime}]
def set_last_error(ql, error_code):
    """Sincronizza LastError sia in Qiling sia all'offset 0x68 del TEB x64 (letto via GS:[0x68] da Go)."""
    ql.os.last_error = error_code
    teb_addr = ql.os.teb.address if (hasattr(ql.os, "teb") and ql.os.teb) else 0x6000000
    ql.mem.write(teb_addr + 0x68, struct.pack("<I", error_code))

def hook_gcenable(ql):
    # Bypass pulito di runtime.gcenable (0x417560) per esecuzione deterministica single-thread
    ret_addr = struct.unpack("<Q", ql.mem.read(ql.arch.regs.rsp, 8))[0]
    ql.arch.regs.rsp += 8
    ql.arch.regs.rip = ret_addr

# Patch in-memory in setup_hooks(): disattiva os/signal.loop (0x53b240) con istruzione RET (0xC3)
ql.hook_address(hook_gcenable, 0x417560)
ql.mem.write(0x53b240, b"\xC3")
\end{lstlisting}

\subsection{Evasione tramite Creazione di un Processo Figlio Nascosto (\texttt{\_HIDDEN\_CHILD})}

\textbf{Funzionamento della Tecnica Anti-Sandbox:} Numerose sandbox dinamiche e analizzatori comportamentali agganciano il monitoraggio unicamente al processo primario generato all'avvio del test: se tale processo termina entro pochi istanti senza compiere azioni ostili, il campione viene classificato prematuramente come benigno mentre la copia secondaria prosegue indisturbata in background. Appena avviato, il RAT in Go generico interroga tramite \texttt{Get\allowbreak{}Environment\allowbreak{}VariableW} la variabile d'ambiente privata \texttt{\_\allowbreak{}HIDDEN\_\allowbreak{}CHILD}. Al lancio iniziale la variabile non \`e definita: la funzione \texttt{main.\allowbreak{}setup\allowbreak{}Hidden\allowbreak{}Process} invoca quindi \texttt{Create\allowbreak{}ProcessW} per istanziare una nuova copia del medesimo eseguibile con i flag di creazione \texttt{0x8080408} (\texttt{DETACHED\_\allowbreak{}PROCESS | CREATE\_\allowbreak{}NO\_\allowbreak{}WINDOW | CREATE\_\allowbreak{}UNICODE\_\allowbreak{}ENVIRONMENT | EXTENDED\_\allowbreak{}STARTUPINFO\_\allowbreak{}PRESENT}), iniettando nel blocco d'ambiente Unicode del figlio la coppia \texttt{\_\allowbreak{}HIDDEN\_\allowbreak{}CHILD=1}, e subito dopo arresta il processo padre tramite \texttt{os.Exit(0)} (all'indirizzo \texttt{0x77de40}).

\textbf{Logica di Risoluzione e Transizione In-Process:} Poich\'e Qiling emula lo spazio di indirizzamento di un singolo processo, consentire l'esecuzione di \texttt{os.Exit(0)} avrebbe troncato immediatamente l'analisi. Si \`e pertanto adottata una strategia chirurgica a doppio stadio (Listato~\ref{lst:gorat_stage1}):
\begin{enumerate}
	\item l'hook su \texttt{Create\allowbreak{}ProcessW} intercetta la richiesta di spawn, decodifica il blocco d'ambiente UTF-16LE confermando la presenza di \texttt{\_\allowbreak{}HIDDEN\_\allowbreak{}CHILD=1}, esporta i metadati nel primo artefatto binario (\texttt{gorat\_\allowbreak{}payload\_\allowbreak{}1.\allowbreak{}bin}), popola una struttura \texttt{PROCESS\_\allowbreak{}INFORMATION} valida (\texttt{PID=1234}, \texttt{TID=5678}) e imposta il flag interno \texttt{ql.\allowbreak{}hidden\_\allowbreak{}child\_\allowbreak{}spawned = True};
	\item immediatamente all'uscita dalla routine di creazione (\texttt{0x77dd94}), dove il codice Go valuta il valore booleano restituito nel registro \texttt{RAX} per decidere se invocare \texttt{os.Exit(0)}, l'hook su indirizzo \texttt{hook\_\allowbreak{}stage\_\allowbreak{}transition} azzera forzatamente \texttt{RAX = 0}. In tal modo il processo padre viene costretto a proseguire senza interruzioni direttamente nella logica dello Stadio 2, assumendo l'identit\`a del processo figlio nascosto (Listato~\ref{lst:gorat_stage1_out}).
\end{enumerate}

\begin{lstlisting}[style=pythoncode, caption={Intercettazione dello spawn nascosto su CreateProcessW e transizione diretta in-process a 0x77dd94.}, label={lst:gorat_stage1}]
@winsdkapi(cc=STDCALL, params={
    "lpApplicationName": WSTRING, "lpCommandLine": WSTRING,
    "lpProcessAttributes": POINTER, "lpThreadAttributes": POINTER,
    "bInheritHandles": BOOL, "dwCreationFlags": DWORD,
    "lpEnvironment": POINTER, "lpCurrentDirectory": WSTRING,
    "lpStartupInfo": POINTER, "lpProcessInformation": POINTER
})
def hook_CreateProcessW(ql, address, params):
    ql.hidden_child_spawned = True
    if params["lpProcessInformation"]:
        fake_pi = struct.pack("<QQII", 0x1111, 0x2222, 1234, 5678)
        ql.mem.write(params["lpProcessInformation"], fake_pi)
    return 1

def hook_stage_transition(ql):
    # A 0x77dd94 RAX=1 causerebbe la chiamata a os.Exit(0) del padre: forziamo RAX=0
    if getattr(ql, "hidden_child_spawned", False):
        ql.arch.regs.rax = 0
\end{lstlisting}


\noindent\begin{minipage}{\linewidth}
\begin{lstlisting}[style=consoleout, caption={Output dell'intercettazione di \_HIDDEN\_CHILD e della transizione in-process dallo Stadio 1 allo Stadio 2.}, label={lst:gorat_stage1_out}]
[DEBUG/ENV] GetEnvironmentVariableW('_HIDDEN_CHILD') -> Non impostato (Attivazione Stage 1: setupHiddenProcess!)
[!] STAGE 1 -> STAGE 2 SPAWN RILEVATO: Chiamata a CreateProcessW
[!] Command Line : "C:\Windows\Documenti\Source\GoRAT\go_rat.exe"
[!] Flags        : 0x8080408 (DETACHED_PROCESS=0x8 | CREATE_NO_WINDOW=0x8000000 | CREATE_UNICODE_ENVIRONMENT=0x400)
[!] Variabili Ambiente Figlio: ['COMPUTERNAME=DESKTOP-GORAT', 'USERNAME=Admin', '_HIDDEN_CHILD=1']
[*] [STAGE TRANSITION @ 0x77dd94] Bypass os.Exit(0) del processo padre!
[*] Transizione diretta in-process dallo Stage 1 (Dropper) allo Stage 2 (Hidden Child Payload _HIDDEN_CHILD=1)!
\end{lstlisting}
\end{minipage}


\subsection[Controllo del Processore tramite COM/OLE e Query WMI]{Controllo del Processore tramite Sottosistema COM/OLE e Query WMI (\texorpdfstring{\texttt{Win32\_\allowbreak{}Processor}}{Win32\_Processor})}

\textbf{Funzionamento della Tecnica Anti-VM:} Entrato nella fase di ricognizione hardware (basata internamente sulla libreria Go \texttt{gopsutil} e sul binding COM \texttt{go-ole}), il RAT in Go generico inizializza il sottosistema COM di Windows (\texttt{Co\allowbreak{}Initialize\allowbreak{}Ex}), risolve il ProgID \texttt{Wbem\allowbreak{}Scripting.\allowbreak{}SWbem\allowbreak{}Locator} tramite \texttt{CLSID\allowbreak{}From\allowbreak{}ProgID}, istanzia l'oggetto con \texttt{Co\allowbreak{}Create\allowbreak{}Instance} ed esegue la query WQL:
\begin{center}
	\texttt{SELECT * FROM Win32\_\allowbreak{}Processor}
\end{center}
Sulle macchine virtuali non opportunamente mascherate tale query restituisce stringhe inequivocabili (quali \texttt{QEMU Virtual CPU}, \texttt{VMware} o \texttt{VirtualBox}) e un numero ridotto di core; in Qiling, invece, l'infrastruttura COM/WMI e il demone \texttt{winmgmt} sono totalmente assenti, causando il dereferenziamento di puntatori non inizializzati e l'interruzione immediata dell'esecuzione.

\textbf{Logica di Risoluzione e Costruzione delle VTable COM in Memoria:} Per superare questo scoglio senza ricorrere a patch distruttive sul binario Go, nello script Python \`e stato costruito da zero in memoria virtuale (a partire dal segmento mappato \texttt{0x70005000}) un intero sottosistema Mock COM/OLE a 64 bit conforme all'ABI Microsoft x64 (Listato~\ref{lst:gorat_wmi}):
\begin{itemize}
	\item \`e stato riscritto l'allocatore OLE delle stringhe \texttt{BSTR} (\texttt{Sys\allowbreak{}Alloc\allowbreak{}String\allowbreak{}Len}, \texttt{SysStringLen}), anteponendo correttamente il prefisso di lunghezza in byte a 32 bit (\texttt{p - 4}) richiesto dalle convenzioni COM e assente nell'implementazione base di Qiling;
	\item agli indirizzi \texttt{0x70005000} (\texttt{SWbem\allowbreak{}Services}/\texttt{SWbem\allowbreak{}Object\allowbreak{}Set}), \texttt{0x70006000} (\texttt{IEnumVARIANT}) e \texttt{0x70007000} (\texttt{SWbemObject} del singolo processore) sono state istanziate vere \textbf{Virtual Method Table (VTable)} a 64 bit i cui slot puntano a \textit{trampoline} contenenti un'istruzione \texttt{RET} (\texttt{0xC3}) agganciati via \texttt{ql.\allowbreak{}hook\_\allowbreak{}address};
	\item quando il codice Go invoca dinamicamente \texttt{IDispatch::\allowbreak{}GetIDs\allowbreak{}Of\allowbreak{}Names} e \texttt{IDispatch::\allowbreak{}Invoke} (incluso il membro speciale \texttt{\_NewEnum} con \texttt{DISPID = -4}, restituito come \texttt{VT\_\allowbreak{}UNKNOWN = 13} per sbloccare l'iteratore \texttt{IEnumVARIANT::\allowbreak{}Next}), i trampoline Python confezionano strutture \texttt{VARIANT} da 24 byte che descrivono una CPU desktop reale \textbf{Intel(R) Core(TM) i7-9700K CPU @ 3.60GHz} (\texttt{GenuineIntel}, 8 core fisici, 8 processori logici, frequenza massima 3600 MHz, \texttt{ProcessorID = BFEBFBFF000906ED}), come attestato dal Listato~\ref{lst:gorat_wmi_out}.
\end{itemize}

\begin{lstlisting}[style=pythoncode, caption={Costruzione delle VTable binarie IDispatch/IEnumVARIANT e spoofing WMI della CPU Intel Core i7-9700K.}, label={lst:gorat_wmi}]
def hook_Invoke_trampoline(ql):
    dispid_raw = ql.arch.regs.rdx & 0xFFFFFFFF
    dispid = dispid_raw - 0x100000000 if dispid_raw >= 0x80000000 else dispid_raw
    p_var = struct.unpack("<Q", ql.mem.read(ql.arch.regs.rsp + 0x38, 8))[0]

    if p_var != 0:
        if dispid == -4:     # DISPID_NEWENUM -> VT_UNKNOWN (13) verso IEnumVARIANT (0x70006000)
            ql.mem.write(p_var, struct.pack("<H6xQ8x", 13, 0x70006000))
        elif dispid == 2002: # Manufacturer -> VT_BSTR (8)
            ql.mem.write(p_var, struct.pack("<H6xQ8x", 8, alloc_bstr(ql, "GenuineIntel")))
        elif dispid == 2003: # Name -> VT_BSTR (8)
            cpu = "Intel(R) Core(TM) i7-9700K CPU @ 3.60GHz"
            ql.mem.write(p_var, struct.pack("<H6xQ8x", 8, alloc_bstr(ql, cpu)))
        elif dispid in (2004, 2005): # NumberOfLogicalProcessors / NumberOfCores -> VT_I4 (3)
            ql.mem.write(p_var, struct.pack("<H6xi12x", 3, 8))
        elif dispid == 2008: # MaxClockSpeed -> VT_I4 (3)
            ql.mem.write(p_var, struct.pack("<H6xi12x", 3, 3600))
    ql.arch.regs.rax = 0  # S_OK
\end{lstlisting}


\noindent\begin{minipage}{\linewidth}
\begin{lstlisting}[style=consoleout, caption={Log dell'interrogazione COM/WMI intercettata e dei parametri hardware restituiti al malware.}, label={lst:gorat_wmi_out}]
[DEBUG/COM] CLSIDFromProgID chiamato per ProgID: 'WbemScripting.SWbemLocator'
[+] WMI ExecQuery Argomento [0]: "SELECT * FROM Win32_Processor"
[+] IDispatch::Invoke: Restituito _NewEnum come VT_UNKNOWN (0x70006000) [UNBLOCK GO REFLECT!]
[+] IEnumVARIANT::Next: Restituito 1 record win32_Processor (VT_DISPATCH: 0x70007000) -> S_OK (0)
[+] IDispatch::Invoke: Restituito Manufacturer = 'GenuineIntel' (VT_BSTR)
[+] IDispatch::Invoke: Restituito Name = 'Intel(R) Core(TM) i7-9700K CPU @ 3.60GHz' (VT_BSTR)
[+] IDispatch::Invoke: Restituito NumberOfLogicalProcessors = 8 (VT_I4)
[+] IDispatch::Invoke: Restituito NumberOfCores = 8 (VT_I4)
[+] IDispatch::Invoke: Restituito MaxClockSpeed = 3600 (VT_I4)
\end{lstlisting}
\end{minipage}


\subsection{Triplice Controllo Temporale: Uptime, \texttt{KUSER\_SHARED\_DATA} e Carico della CPU}

\textbf{Funzionamento della Tecnica Anti-Sandbox:} Il RAT in Go generico intreccia tre verifiche temporali complementari che mettono in crisi sia le sandbox appena avviate sia gli emulatori binari:
\begin{enumerate}
	\item interroga \texttt{Get\allowbreak{}Tick\allowbreak{}Count64} per calcolare l'\textit{uptime} complessivo della macchina, scartando gli ambienti virtuali accesi da poche decine di secondi;
	\item invoca due volte \texttt{Nt\allowbreak{}Query\allowbreak{}System\allowbreak{}Information} con classe \texttt{System\allowbreak{}Processor\allowbreak{}Performance\allowbreak{}Information} (\texttt{8}), interponendo una pausa esatta di un secondo (\texttt{time.\allowbreak{}Sleep(1 * time.\allowbreak{}Second)}) per calcolare la variazione dei contatori \texttt{IdleTime}, \texttt{KernelTime} e \texttt{UserTime} degli 8 core e ricavare la percentuale effettiva di utilizzo della CPU;
	\item durante \texttt{time.Sleep}, il runtime Go attende su una porta di completamento IOCP (\texttt{Get\allowbreak{}Queued\allowbreak{}Completion\allowbreak{}Status\allowbreak{}Ex}) e verifica l'effettivo trascorrere del tempo leggendo direttamente in memoria utente i campi a 64 bit \texttt{Interrupt\allowbreak{}Time} (\texttt{0x7ffe0008}) e \texttt{SystemTime} (\texttt{0x7ffe0014}) nella pagina condivisa del kernel \texttt{KUSER\_\allowbreak{}SHARED\_\allowbreak{}DATA}. Poich\'e in Qiling tale pagina non viene aggiornata da alcun interrupt timer hardware, il runtime Go rileva tempo trascorso pari a zero e richiama indefinitamente l'attesa, bloccando l'analisi in un ciclo infinito.
\end{enumerate}

\textbf{Logica di Risoluzione e Implementazione:} La routine \texttt{update\_\allowbreak{}kuser\_\allowbreak{}time} sviluppata nello script fa avanzare programmaticamente i contatori a unit\`a di 100 nanosecondi in \texttt{0x7ffe0008}, \texttt{0x7ffe0014} e \texttt{0x7ffe0320} ad ogni invocazione di pausa o timeout IOCP, risolvendo istantaneamente il deadlock del runtime Go. Contestualmente, \texttt{Get\allowbreak{}Tick\allowbreak{}Count64} restituisce un sistema operativo attivo da oltre un'ora (\texttt{3604 s}), mentre l'hook su \texttt{Nt\allowbreak{}Query\allowbreak{}System\allowbreak{}Information} popola per ciascuno degli 8 core strutture da 48 byte con contatori progressivamente crescenti tra la prima e la seconda lettura, calibrati affinch\'e la formula differenziale del malware calcoli un carico fisiologico della CPU pari esattamente al \textbf{24,69\%} (Listati~\ref{lst:gorat_time} e \ref{lst:gorat_time_out}).

\begin{lstlisting}[style=pythoncode, caption={Avanzamento di KUSER\_SHARED\_DATA, falsificazione dell'uptime (>1 ora) e simulazione del carico CPU.}, label={lst:gorat_time}]
def update_kuser_time(ql, increment_ms=1000):
    curr_time = struct.unpack("<Q", ql.mem.read(0x7ffe0008, 8))[0]
    curr_time += int(increment_ms) * 10000  # Conversione ms -> unita' da 100 ns
    ql.mem.write(0x7ffe0008, struct.pack("<Q", curr_time))
    ql.mem.write(0x7ffe0014, struct.pack("<Q", curr_time))
    ql.mem.write(0x7ffe0320, struct.pack("<I", int(curr_time // 10000) & 0xFFFFFFFF))

@winsdkapi(cc=STDCALL, params={})
def hook_GetTickCount64(ql, address, params):
    global tempo_simulato
    tempo_simulato += 250
    update_kuser_time(ql, 250)
    return 3600000 + tempo_simulato  # Oltre 1 ora di accensione simulata

@winsdkapi(cc=STDCALL, params={
    "SystemInformationClass": ULONG, "SystemInformation": POINTER,
    "SystemInformationLength": ULONG, "ReturnLength": POINTER
})
def hook_NtQuerySystemInformation(ql, address, params):
    global nt_cpu_call_count
    if params["SystemInformationClass"] == 8:  # SystemProcessorPerformanceInformation
        nt_cpu_call_count += 1
        s = nt_cpu_call_count
        idle   = 100000000 + s * 7550000
        kernel = 150000000 + s * 9000000  # Su Windows KernelTime include IdleTime
        user   = 30000000  + s * 1000000
        entry = struct.pack("<qqqqqQ", idle, kernel, user, 2000000, 1000000, 5000)
        ql.mem.write(params["SystemInformation"], entry * 8)
        return 0
\end{lstlisting}


\noindent\begin{minipage}{\linewidth}
\begin{lstlisting}[style=consoleout, caption={Superamento del campionamento CPU su 8 core, avanzamento IOCP/KUSER\_SHARED\_DATA e verifica dell'uptime.}, label={lst:gorat_time_out}]
[DEBUG/SYSINFO] NtQuerySystemInformation(SystemProcessorPerformanceInformation): Popolati 8 core (campione #1)
[DEBUG/IOCP] GetQueuedCompletionStatusEx: Avanzato orologio KUSER_SHARED_DATA di +1000ms (Pausa 1s superata)
[DEBUG/SYSINFO] NtQuerySystemInformation(SystemProcessorPerformanceInformation): Popolati 8 core (campione #2)
[DEBUG/SYSINFO] GetTickCount64 chiamato -> Restituito uptime=3604500 ms (3604 s)
\end{lstlisting}
\end{minipage}


\subsection{Verifica delle Risorse Fisiche: Memoria RAM, Disco NTFS, Traffico di Rete e Sistema Operativo}

\textbf{Funzionamento della Tecnica Anti-VM:} Le macchine virtuali usa-e-getta dedicate all'analisi automatica vengono tipicamente allocate con profili hardware minimali: 2--4 GB di RAM, dischi virtuali inferiori a 60 GB e contatori di byte trasmessi/ricevuti sull'interfaccia di rete prossimi allo zero. Il RAT in Go generico sonda sistematicamente ciascuno di questi indicatori invocando \texttt{Global\allowbreak{}Memory\allowbreak{}Status\allowbreak{}Ex}, \texttt{Get\allowbreak{}Logical\allowbreak{}Drive\allowbreak{}StringsW}, \texttt{Get\allowbreak{}Volume\allowbreak{}InformationW}, \texttt{Get\allowbreak{}Disk\allowbreak{}Free\allowbreak{}Space\allowbreak{}ExW}, \texttt{Get\allowbreak{}Adapters\allowbreak{}Addresses}, \texttt{GetIfEntry2} e \texttt{Rtl\allowbreak{}Get\allowbreak{}Version} (integrata dalla lettura delle chiavi \texttt{Current\allowbreak{}Version} del Registro di Windows).

\textbf{Logica di Risoluzione e Implementazione:} Agganciando l'intera famiglia di API hardware e di rete, l'harness presenta al trojan il profilo organico di una postazione aziendale intensamente utilizzata (Listati~\ref{lst:gorat_hw} e \ref{lst:gorat_hw_out}): \textbf{16 GB di RAM fisica} (6 GB occupati, 10 GB disponibili), volume \texttt{C:\textbackslash\allowbreak{}} formattato in \textbf{NTFS da 500 GB} (180 GB in uso, 320 GB liberi), scheda \texttt{Ethernet0} con IP locale \texttt{192.168.1.100}, MAC \texttt{00:0c:29:12:34:56} e contatori binari \texttt{MIB\_\allowbreak{}IF\_\allowbreak{}ROW2} popolati agli offset \texttt{+0x4b8} e \texttt{+0x500} con \textbf{154,8 MB ricevuti} e \textbf{48,2 MB inviati}, su sistema operativo \textbf{Microsoft Windows 10 Pro Build 19045.3803}.

\begin{lstlisting}[style=pythoncode, caption={Simulazione di 16 GB di RAM e traffico storico di rete nella struttura MIB\_IF\_ROW2.}, label={lst:gorat_hw}]
@winsdkapi(cc=STDCALL, params={"lpBuffer": POINTER})
def hook_GlobalMemoryStatusEx(ql, address, params):
    total_ram = 16 * (1024**3)  # 16 GB totali
    free_ram  = 10 * (1024**3)  # 10 GB disponibili (6 GB in uso)
    data = struct.pack(
        "<IIQQQQQQQ", 64, 37, total_ram, free_ram,
        32 * (1024**3), 24 * (1024**3), 0x7FFFFFFE0000, 0x7FFFFFF00000, 0
    )
    ql.mem.write(params["lpBuffer"], data)
    return 1

@winsdkapi(cc=STDCALL, params={"pIfRow": POINTER})
def hook_GetIfEntry2(ql, address, params):
    p_row = params["pIfRow"]
    if p_row:
        ql.mem.write(p_row + 0x4B8, struct.pack("<Q", 154829104))  # InOctets (~154.8 MB)
        ql.mem.write(p_row + 0x500, struct.pack("<Q", 48291042))   # OutOctets (~48.2 MB)
    return 0
\end{lstlisting}


\noindent\begin{minipage}{\linewidth}
\begin{lstlisting}[style=consoleout, caption={Profilo hardware, disco, rete e sistema operativo acquisito dal RAT in Go generico.}, label={lst:gorat_hw_out}]
[DEBUG/SYSINFO] GlobalMemoryStatusEx: Total=16 GB, Used=6 GB, Free=10 GB
[DEBUG/SYSINFO] GetVolumeInformationW: Volume='Windows', FileSystem='NTFS'
[DEBUG/SYSINFO] GetDiskFreeSpaceExW: Total=500 GB, Used=180 GB, Free=320 GB
[+] GetAdaptersAddresses: Popolato adattatore virtuale 'Ethernet0' (IP: 192.168.1.100, MAC: 00:0c:29:12:34:56)
[DEBUG/SYSINFO] GetIfEntry2: Popolati contatori traffico MIB_IF_ROW2 (InOctets=154.8 MB, OutOctets=48.2 MB)
[DEBUG/SYSINFO] RtlGetVersion: Windows 10.0 Build 19045 (Workstation)
\end{lstlisting}
\end{minipage}


\subsection{Verifica della Connettivit\`a Internet, Geolocalizzazione e Scansione delle Porte Locali}

\textbf{Funzionamento della Tecnica:} Prima di esporsi contattando il server C2, il malware accerta l'effettiva disponibilit\`a di accesso a Internet interrogando due servizi legittimi esterni --- \texttt{api.\allowbreak{}ipify.\allowbreak{}org} per determinare l'indirizzo IP pubblico WAN e \texttt{ip-api.com} per ricavare Nazione e Citt\`a della vittima (\textit{geo-fencing}) --- ed esegue infine una scansione TCP locale (\texttt{sysinfo.\allowbreak{}is\allowbreak{}Port\allowbreak{}Open}) su \textbf{23 porte standard} di \texttt{127.0.0.1} tramite il puntatore di estensione Winsock \texttt{ConnectEx} ottenuto via \texttt{WSAIoctl} (\texttt{SIO\_\allowbreak{}GET\_\allowbreak{}EXTENSION\_\allowbreak{}FUNCTION\_\allowbreak{}POINTER = 0xC8000006}).

\textbf{Logica di Risoluzione e Implementazione:} Poich\'e la query verso \texttt{api.\allowbreak{}ipify.\allowbreak{}org} era codificata originariamente con schema \texttt{https://}, all'avvio dell'emulazione lo script effettua una patch trasparente in memoria di pari lunghezza (33 byte) all'indirizzo \texttt{0x842cb5}, convertendo \texttt{https://\allowbreak{}api.\allowbreak{}ipify.\allowbreak{}org?format=json} in \texttt{http://\allowbreak{}api.\allowbreak{}ipify.\allowbreak{}org/\allowbreak{}?format=json} cos\`i da intercettare la richiesta HTTP in chiaro senza dover negoziare sessioni TLS. Gli hook su \texttt{GetAddrInfoW}, \texttt{WSASend} e \texttt{WSARecv} (sincronizzati con la coda di eventi asincroni IOCP del netpoller Go) forniscono quindi due risposte HTTP JSON sintetiche che attribuiscono all'host l'IP \texttt{198.51.100.42} localizzato a \textbf{New York (United States)}. In parallelo, il trampoline registrato a \texttt{0x70003000} per \texttt{ConnectEx} segnala come aperte le porte locali \texttt{80}, \texttt{443}, \texttt{3389} e \texttt{8080}, rifiutando istantaneamente con errore sincrono \texttt{WSAECONNREFUSED} (\texttt{10061}) le altre 19 porte per azzerare i timeout di rete (Listati~\ref{lst:gorat_net} e \ref{lst:gorat_net_out}).

\begin{lstlisting}[style=pythoncode, caption={Simulazione HTTP per IP pubblico/geolocalizzazione e filtro porte nel trampoline ConnectEx.}, label={lst:gorat_net}]
# Risposte HTTP simulate per api.ipify.org e ip-api.com
simulated_ipify_response = (
    b"HTTP/1.1 200 OK\r\nContent-Type: application/json\r\nContent-Length: 23\r\n"
    b"Connection: close\r\n\r\n{\"ip\":\"198.51.100.42\"}"
)
simulated_ipapi_body = (
    b'{"status":"success","country":"United States","city":"New York","query":"198.51.100.42"}'
)

# Nel trampoline ConnectEx (0x70003000): apre 80, 443, 3389, 8080 su 127.0.0.1 e rifiuta le altre
if ip == "127.0.0.1" and port not in (80, 443, 3389, 8080):
    set_last_error(ql, 10061)  # WSAECONNREFUSED sincrono immediato
    ql.arch.regs.rax = 0
    return
\end{lstlisting}


\noindent\begin{minipage}{\linewidth}
\begin{lstlisting}[style=consoleout, caption={Risoluzione DNS simulata, risposte di geolocalizzazione e scansione delle porte locali.}, label={lst:gorat_net_out}]
[DEBUG/DNS] GetAddrInfoW: Risolto 'api.ipify.org' -> 104.26.12.159:80
[DEBUG/HTTP] WSASend -> Richiesta HTTP GET a api.ipify.org (99 bytes)
[DEBUG/DNS] GetAddrInfoW: Risolto 'ip-api.com' -> 208.95.112.1:80
[DEBUG/HTTP] WSASend -> Richiesta HTTP GET Geolocalizzazione a ip-api.com (104 bytes)
[DEBUG/RECON] Scansione porta locale (sysinfo.isPortOpen): 127.0.0.1:80   -> APERTA
[DEBUG/RECON] Scansione porta locale (sysinfo.isPortOpen): 127.0.0.1:443  -> APERTA
[DEBUG/RECON] Scansione porta locale (sysinfo.isPortOpen): 127.0.0.1:3389 -> APERTA
[DEBUG/RECON] Scansione porta locale (sysinfo.isPortOpen): 127.0.0.1:8080 -> APERTA
\end{lstlisting}
\end{minipage}


\subsection{Simulazione del Server C2 WebSocket (RFC 6455) e Ciclo Interattivo Protocol Buffers}

Superate tutte le verifiche ambientali, il RAT in Go generico genera l'identificativo persistente della vittima \texttt{client-\allowbreak{}DESKTOP-\allowbreak{}GORAT-\allowbreak{}000c29-\allowbreak{}51616}, salvandolo in \texttt{App\allowbreak{}Data\textbackslash\allowbreak{}Roaming\textbackslash\allowbreak{}Remote\allowbreak{}Admin\textbackslash\allowbreak{}client\_\allowbreak{}id}, e apre una connessione TCP verso il server di comando \texttt{54.233.190.22} sulla porta \texttt{8080} iniziando l'\textit{handshake} \textbf{WebSocket} sull'endpoint \texttt{ws://\allowbreak{}54.233.190.22:8080/\allowbreak{}ws}.

Per esplorare integralmente le capacit\`a post-sfruttamento del RAT senza esporre alcuna rete reale, direttamente all'interno di \texttt{hook\_\allowbreak{}WSA\allowbreak{}Send} \`e stato implementato un \textbf{Server C2 WebSocket + Protocol Buffers completo} (Listato~\ref{lst:gorat_c2}):
\begin{enumerate}
	\item intercettata la richiesta HTTP \texttt{Upgrade: websocket}, lo script estrae l'header \texttt{Sec-\allowbreak{}Web\allowbreak{}Socket-\allowbreak{}Key}, calcola dinamicamente il digest SHA-1 concatenato alla costante RFC 6455 (\texttt{258EAFA5-\allowbreak{}E914-\allowbreak{}47DA-\allowbreak{}95CA-\allowbreak{}C5AB0DC85B11}) codificandolo in Base64 e risponde con \texttt{HTTP/\allowbreak{}1.1 101 Switching Protocols};
	\item decodificata la maschera XOR a 4 byte del primo frame binario (\texttt{opcode = 2}) inviato dal trojan, un parser \textit{wire-format} Protocol Buffers custom ricostruisce il messaggio \texttt{proto.\allowbreak{}Client\allowbreak{}Info} contenente l'intera scheda tecnica della macchina appena profilata;
	\item in risposta, il server simulato serializza in formato binario \texttt{proto.\allowbreak{}Command} e incapsula in frame WebSocket server-to-client quattro ordini operativi consecutivi per collaudare tutte le primitive del RAT (oltre al supporto per l'esecuzione di comandi da terminale): \texttt{pwd} (lettura directory corrente), \texttt{upload\_file} (scrittura di un eseguibile inviato dall'attaccante in \texttt{C:\textbackslash\allowbreak{}Windows\textbackslash\allowbreak{}Temp\textbackslash\allowbreak{}c2\_\allowbreak{}dropped\_\allowbreak{}module.\allowbreak{}bin}), \texttt{download\_\allowbreak{}file} (esfiltrazione del file \texttt{client\_id} dalla vittima) e \texttt{delete\_file} (rimozione chirurgica del modulo precedentemente caricato);
	\item ricevuta la quarta conferma, il C2 simulato invia un \textit{Close Frame} WebSocket (\texttt{0x88 0x02 0x03 0xe9}, codice \texttt{1001 Going Away}), consentendo al malware di chiudere il socket e terminare l'esecuzione in modo totalmente pulito.
\end{enumerate}

\begin{lstlisting}[style=pythoncode, caption={Handshake WebSocket RFC 6455, demasking dei frame binari e iniezione di comandi Protocol Buffers.}, label={lst:gorat_c2}]
def encode_proto_command(cmd_id, command, args=None, file_data=None):
    """Serializza un comando remoto proto.Command nel wire-format binario Protocol Buffers."""
    buf = bytearray()
    buf.extend(encode_proto_bytes(1, cmd_id.encode("utf-8")))
    buf.extend(encode_proto_bytes(2, command.encode("utf-8")))
    for a in (args or []):
        buf.extend(encode_proto_bytes(3, a.encode("utf-8")))
    if file_data:
        buf.extend(encode_proto_bytes(4, file_data))
    return bytes(buf)

# Calcolo dinamico di Sec-WebSocket-Accept secondo RFC 6455 in hook_WSASend:
accept_raw = hashlib.sha1((ws_key + "258EAFA5-E914-47DA-95CA-C5AB0DC85B11").encode("ascii")).digest()
accept_key = base64.b64encode(accept_raw).decode("ascii")
ws_resp = (
    f"HTTP/1.1 101 Switching Protocols\r\nUpgrade: websocket\r\n"
    f"Connection: Upgrade\r\nSec-WebSocket-Accept: {accept_key}\r\n\r\n"
).encode("ascii")
\end{lstlisting}

La Tabella~\ref{tab:gorat_payloads} riassume i sette artefatti binari salvati automaticamente durante la sessione, mentre il Listato~\ref{lst:gorat_json} riporta il report JSON integrale decodificato dallo scambio Protobuf (\texttt{gorat\_\allowbreak{}c2\_\allowbreak{}full\_\allowbreak{}report.\allowbreak{}json}).

\begin{table}[!htbp]
	\centering
	\footnotesize
	\setlength{\tabcolsep}{5pt}
	\renewcommand{\arraystretch}{1.32}
	\begin{tabularx}{\textwidth}{|l|r|Y|}
		\hline
		\textbf{File Estratto} & \textbf{Dim.} & \textbf{Contenuto Tecnico Catturato} \\
		\hline
		\texttt{gorat\_\allowbreak{}payload\_\allowbreak{}1.\allowbreak{}bin} & 217 B & Metadati e blocco ambiente di \texttt{Create\allowbreak{}ProcessW} nello Stadio 1 (\texttt{\_HIDDEN\_\allowbreak{}CHILD=1}) \\
		\hline
		\texttt{gorat\_\allowbreak{}payload\_\allowbreak{}2.\allowbreak{}bin} & 412 B & Messaggio Protobuf iniziale \texttt{proto.\allowbreak{}ClientInfo} con l'intero profilo hardware/rete della vittima \\
		\hline
		\texttt{gorat\_\allowbreak{}payload\_\allowbreak{}3.\allowbreak{}bin} & 91 B & Risposta Protobuf al comando \texttt{pwd} (\texttt{C:\textbackslash\allowbreak{}Windows\textbackslash\allowbreak{}Documenti\textbackslash\allowbreak{}Source\textbackslash\allowbreak{}GoRAT}) \\
		\hline
		\texttt{gorat\_\allowbreak{}payload\_\allowbreak{}4.\allowbreak{}bin} & 52 B & Modulo PE di test ricevuto dal C2 e scritto su disco tramite il comando \texttt{upload\_\allowbreak{}file} \\
		\hline
		\texttt{gorat\_\allowbreak{}payload\_\allowbreak{}5.\allowbreak{}bin} & 81 B & Risposta Protobuf di conferma scrittura file (\texttt{"Successfully wrote file"}) \\
		\hline
		\texttt{gorat\_\allowbreak{}payload\_\allowbreak{}6.\allowbreak{}bin} & 120 B & Risposta Protobuf al comando \texttt{download\_\allowbreak{}file} contenente i byte esfiltrati di \texttt{client\_\allowbreak{}id} \\
		\hline
		\texttt{gorat\_\allowbreak{}payload\_\allowbreak{}7.\allowbreak{}bin} & 93 B & Risposta Protobuf di conferma cancellazione file tramite \texttt{delete\_\allowbreak{}file} \\
		\hline
	\end{tabularx}
	\caption{Artefatti binari catturati ed esportati automaticamente durante l'emulazione completa del RAT in Go generico.}
	\label{tab:gorat_payloads}
\end{table}

\begin{lstlisting}[style=consoleout, caption={Contenuto integrale del report C2 decodificato automaticamente da Protocol Buffers (gorat\_c2\_full\_report.json).}, label={lst:gorat_json}]
{
  "client_info": {
    "id": "client-DESKTOP-GORAT-000c29-51616",
    "hostname": "DESKTOP-GORAT",
    "os": "windows",
    "ip_address": "198.51.100.42",
    "country": "United States",
    "city": "New York",
    "system_info": {
      "cpu_cores": 1,
      "cpu_usage_percent": 24.69,
      "cpu_model": "Intel(R) Core(TM) i7-9700K CPU @ 3.60GHz",
      "total_memory_gb": 16.0,
      "used_memory_gb": 6.0,
      "free_memory_gb": 10.0,
      "disks": [
        {
          "device": "C:",
          "mount_point": "C:",
          "total_space_gb": 500.0,
          "used_space_gb": 180.0,
          "free_space_gb": 320.0,
          "filesystem_type": "NTFS"
        }
      ],
      "network_interfaces": [
        {
          "name": "Ethernet0",
          "ip_addresses": ["192.168.1.100/0"],
          "mac_address": "00:0c:29:12:34:56",
          "bytes_sent": 48291042,
          "bytes_received": 154829104
        }
      ],
      "os_version": "Microsoft Windows 10 Pro  10.0.19045.3803 Build 19045.3803",
      "kernel_version": "10.0.19045.3803 Build 19045.3803",
      "uptime_seconds": 3604,
      "open_ports": [80, 443, 3389, 8080]
    }
  },
  "c2_commands_responses": [
    {"id": "cmd-001-pwd", "output": "C:\\Windows\\Documenti\\Source\\GoRAT", "exit_code": 0, "completed": true},
    {"id": "cmd-002-upload", "output": "Successfully wrote file", "exit_code": 0, "completed": true},
    {"id": "cmd-003-download", "output": "Successfully read file", "exit_code": 0, "completed": true,
     "file_data_len": 33, "file_data_preview": "client-DESKTOP-GORAT-000c29-51616"},
    {"id": "cmd-004-delete", "output": "Successfully deleted file/directory", "exit_code": 0, "completed": true}
  ]
}
\end{lstlisting}

% =====================================================================
% SEZIONE 7
% =====================================================================
\section{Caso di Studio 3: Unpacking Dinamico UPX, Simulazione CryptoAPI e Neutralizzazione DoS nel Malware Astaroth}

Il terzo caso di studio prende in esame \texttt{Astaroth.exe} (famiglia \textit{Astaroth / Guildma}), un eseguibile Windows PE32 a 32 bit compilato in C mediante toolchain \textit{MinGW GCC 13} e compresso con \textbf{UPX 4.22}. Durante le prime prove di esecuzione non strumentata, il binario mandava immediatamente in crash l'emulatore impedendo qualsiasi osservazione comportamentale. Attraverso lo sviluppo dell'harness dedicato \texttt{astaroth\_\allowbreak{}upx\_\allowbreak{}unpacker\_\allowbreak{}analyzer.\allowbreak{}py} \`e stato possibile non solo lasciare che il malware si decompattasse autonomamente in memoria, ma anche neutralizzare sei distinti presidi di offuscamento, anti-analisi e sabotaggio delle risorse, completando l'esecuzione fino alla naturale terminazione del processo.

\subsection{Compressione tramite Packer UPX (\texttt{UPX 4.22}) ed Estrazione dalla Memoria RAM}

\textbf{Funzionamento della Tecnica:} All'ispezione statica su disco, l'intestazione PE del campione non presenta le consuete sezioni \texttt{.text}, \texttt{.rdata} e \texttt{.data}, bens\`i tre segmenti denominati \texttt{UPX0}, \texttt{UPX1} e \texttt{UPX2}; la tabella delle importazioni espone appena sette API generiche utilizzate dallo \textit{stub} di decompressione (situato all'Entry Point \texttt{0x00413360}), mentre tutte le stringhe operative e le funzioni malevole risultano compresse e illeggibili.

\textbf{Logica di Risoluzione e In-Memory Dumping:} Sfruttando l'emulazione istruzione per istruzione della CPU x86 fornita da Qiling, il binario \`e stato avviato direttamente dallo stub UPX (\texttt{0x00413360}), lasciando che la routine nativa ricostruisse in autonomia l'immagine PE originale in RAM e risolvesse dinamicamente le proprie importazioni. Posizionando tramite \texttt{ql.\allowbreak{}hook\_\allowbreak{}address} un intercettore chirurgico all'ingresso della funzione \texttt{main()} del codice decompattato (\texttt{0x00401a18}, subito a valle dell'\textit{Original Entry Point} CRT \texttt{0x004014c0}), lo script cattura istantaneamente \texttt{0x15000} byte (86.016 byte) a partire da \texttt{IMAGE\_\allowbreak{}BASE = 0x00400000} e li riversa in \texttt{astaroth\_\allowbreak{}unpacked\_\allowbreak{}dump.\allowbreak{}bin}, esponendo l'intero codice disassemblabile del malware (Listato~\ref{lst:astaroth_upx}).

\begin{lstlisting}[style=pythoncode, caption={Decompattamento automatico UPX in Qiling e dump dell'immagine PE ricostruita all'ingresso di main().}, label={lst:astaroth_upx}]
IMAGE_BASE = 0x00400000
UNPACKED_DUMP_SIZE = 0x15000
ADDR_UNPACKED_MAIN = 0x00401A18

def hook_unpacked_main(ql):
    print(f"[+] Decompattamento UPX completato! Raggiunto main() a 0x{ADDR_UNPACKED_MAIN:08x}.")
    dump_path = os.path.join(_output_dir, "astaroth_unpacked_dump.bin")
    mem = ql.mem.read(IMAGE_BASE, UNPACKED_DUMP_SIZE)
    with open(dump_path, "wb") as f:
        f.write(mem)
    print(f"[+] Dump della memoria PE decompattata salvato in: {dump_path}")

ql.hook_address(hook_unpacked_main, ADDR_UNPACKED_MAIN)
\end{lstlisting}

\subsection{Offuscamento Stringhe via XOR (\texttt{0x5A}) e Trappola Anti-Emulatore su CryptoAPI (\texttt{AES-256})}

\textbf{Funzionamento della Tecnica:} All'ingresso di \texttt{main()}, Astaroth alloca sullo stack la sequenza di 11 byte offuscati \texttt{17 35 36 35 34 7a 16 3b 3b 38 3f} e vi applica un ciclo XOR con chiave a singolo byte \texttt{0x5A}, ricostruendo solo temporaneamente in RAM la frase segreta \texttt{"Molon Labe"} (celebre motto spartano impiegato dall'autore come firma interna). Immediatamente dopo, il programma invoca una catena di cinque API del sottosistema crittografico Microsoft (\texttt{ADVAPI32.dll}): \texttt{Crypt\allowbreak{}Acquire\allowbreak{}ContextA} (richiedendo il provider \texttt{PROV\_\allowbreak{}RSA\_\allowbreak{}AES = 0x18}), \texttt{CryptGenKey} (con algoritmo \texttt{CALG\_\allowbreak{}AES\_\allowbreak{}256 = 0x6610}), \texttt{CryptEncrypt}, \texttt{Crypt\allowbreak{}Destroy\allowbreak{}Key} e \texttt{Crypt\allowbreak{}Release\allowbreak{}Context}. Tale sequenza funge da efficacissima \textbf{trappola Anti-Emulatore}: poich\'e Qiling non implementa nativamente \texttt{Crypt\allowbreak{}Acquire\allowbreak{}ContextA} n\'e l'infrastruttura interna di \texttt{cryptsp.dll}, il tentativo di eseguire il codice grezzo della DLL portava a un crash fatale per accesso invalido alla memoria (\texttt{Invalid memory fetch} a \texttt{0x00012c1e}), arrestando l'analisi prima di qualunque attivit\`a su disco o rete.

\textbf{Logica di Risoluzione ed Estrazione Just-in-Time:} Nell'harness sono stati implementati tutti e cinque gli hook per le funzioni CryptoAPI restituendo handle sintetici coerenti (\texttt{0xcafe0001} per il provider e \texttt{0xcafe0002} per la chiave AES-256). In particolare, intercettando \texttt{CryptEncrypt} nell'esatto istante in cui riceve il puntatore allo stack \texttt{pbData}, lo script legge i byte appena decifrati dal ciclo XOR estraendo in chiaro \texttt{"Molon Labe"}, applica una trasformazione simulata in-place sul buffer e restituisce codice di successo \texttt{1} (Listati~\ref{lst:astaroth_crypto} e \ref{lst:astaroth_crypto_out}).

\begin{lstlisting}[style=pythoncode, caption={Simulazione del sottosistema Microsoft CryptoAPI ed estrazione in chiaro del buffer decifrato via XOR.}, label={lst:astaroth_crypto}]
@winsdkapi(cc=STDCALL, params={
    "phProv": POINTER, "szContainer": LPCSTR, "szProvider": LPCSTR,
    "dwProvType": DWORD, "dwFlags": DWORD
})
def hook_CryptAcquireContextA(ql, address, params):
    if params["phProv"]:
        ql.mem.write(params["phProv"], struct.pack("<I", 0xCAFE0001))
    return 1

@winsdkapi(cc=STDCALL, params={
    "hKey": HANDLE, "hHash": HANDLE, "Final": INT, "dwFlags": DWORD,
    "pbData": POINTER, "pdwDataLen": POINTER, "dwBufLen": DWORD
})
def hook_CryptEncrypt(ql, address, params):
    pbData = params["pbData"]
    data_len = struct.unpack("<I", ql.mem.read(params["pdwDataLen"], 4))[0]
    raw = bytes(ql.mem.read(pbData, data_len))
    plaintext_str = raw.split(b"\x00")[0].decode("utf-8", errors="replace")
    ql.mem.write(pbData, bytes((b ^ 0xA5) for b in raw))  # Cifratura simulata in-place
    telemetry["decrypted_strings"].append(plaintext_str)
    return 1
\end{lstlisting}


\noindent\begin{minipage}{\linewidth}
\begin{lstlisting}[style=consoleout, caption={Traccia delle chiamate CryptoAPI intercettate e recupero della firma 'Molon Labe'.}, label={lst:astaroth_crypto_out}]
[CRYPTO] CryptAcquireContextA(dwProvType=0x18 [PROV_RSA_AES], dwFlags=0xf0000000) -> hProv=0xcafe0001
[CRYPTO] CryptGenKey(Algid=0x6610 [CALG_AES_256], dwFlags=0x1) -> hKey=0xcafe0002
[CRYPTO] CryptEncrypt(pbData=0xffffcf7d, len=11, bufLen=27, plaintext='Molon Labe')
[CRYPTO] CryptDestroyKey(hKey=0xcafe0002)
[CRYPTO] CryptReleaseContext(hProv=0xcafe0001)
\end{lstlisting}
\end{minipage}


\subsection{Rilevamento Debugger (\texttt{IsDebuggerPresent}) ed Elevazione Privilegi UAC (\texttt{runas})}

\textbf{Funzionamento della Tecnica e Bypass:} Superato il blocco crittografico, all'indirizzo \texttt{0x00401b43} il malware invoca \texttt{Is\allowbreak{}Debugger\allowbreak{}Present()} per ispezionare il flag \texttt{Being\allowbreak{}Debugged} nel PEB del processo e, subito dopo (\texttt{0x00401b89}), richiama \texttt{Shell\allowbreak{}ExecuteA} passando il verbo \texttt{"runas"} sul proprio eseguibile al fine di ottenere i privilegi di Amministratore tramite prompt UAC. L'harness neutralizza entrambe le manovre forzando \texttt{Is\allowbreak{}Debugger\allowbreak{}Present} a restituire \texttt{0} e facendo restituire a \texttt{Shell\allowbreak{}ExecuteA} il valore \texttt{42} (nelle specifiche Win32 ogni codice $> 32$ attesta la riuscita dell'avvio), impedendo la generazione di un processo esterno non controllato (Listato~\ref{lst:astaroth_debug}).

\begin{lstlisting}[style=pythoncode, caption={Bypass del controllo anti-debug e simulazione del successo dell'elevazione privilegi UAC.}, label={lst:astaroth_debug}]
@winsdkapi(cc=STDCALL, params={})
def hook_IsDebuggerPresent(ql, address, params):
    return 0  # Nessun debugger rilevato

@winsdkapi(cc=STDCALL, params={
    "hwnd": HANDLE, "lpOperation": LPCSTR, "lpFile": LPCSTR,
    "lpParameters": LPCSTR, "lpDirectory": LPCSTR, "nShowCmd": INT
})
def hook_ShellExecuteA(ql, address, params):
    print(f"[PRIV-ESC] ShellExecuteA(op='{params['lpOperation']}', file='{params['lpFile']}')")
    return 42  # Valore > 32 segnala pieno successo al chiamante
\end{lstlisting}

\subsection{Occultamento sul File System (\texttt{5kidRo0t.exe}) e Doppia Persistenza nel Registro}

Nella subroutine \texttt{0x0040150d}, il campione interroga \texttt{SH\allowbreak{}Get\allowbreak{}Folder\allowbreak{}PathA} con parametro \texttt{CSIDL\_\allowbreak{}APPDATA} (\texttt{0x1a}), concatena al percorso restituito il nome offuscato \texttt{\textbackslash\allowbreak{}5kid\allowbreak{}Ro0t.\allowbreak{}exe} e vi si clona tramite \texttt{CopyFileA}. Successivamente apre in sequenza la chiave \texttt{Software\textbackslash\allowbreak{}Microsoft\textbackslash\allowbreak{}Windows\textbackslash\allowbreak{}Current\allowbreak{}Version\textbackslash\allowbreak{}Run} sia nel ramo utente \texttt{HKEY\_\allowbreak{}CURRENT\_\allowbreak{}USER} (che garantisce la persistenza anche senza privilegi elevati) sia nel ramo di sistema \texttt{HKEY\_\allowbreak{}LOCAL\_\allowbreak{}MACHINE} (sfruttando l'eventuale elevazione UAC), registrando in entrambi il valore \texttt{5kidRo0t}. Infine invoca \texttt{Set\allowbreak{}File\allowbreak{}AttributesA} con maschera \texttt{0x6} (\texttt{FILE\_\allowbreak{}ATTRIBUTE\_\allowbreak{}HIDDEN | FILE\_\allowbreak{}ATTRIBUTE\_\allowbreak{}SYSTEM}), trasformando la copia in un file protetto di sistema invisibile all'esplorazione standard di Windows (Listato~\ref{lst:astaroth_persist_out}).


\noindent\begin{minipage}{\linewidth}
\begin{lstlisting}[style=consoleout, caption={Tracciamento completo dell'occultamento su file system e della doppia persistenza nel Registro.}, label={lst:astaroth_persist_out}]
[FILESYSTEM] SHGetFolderPathA(csidl=0x1a [CSIDL_APPDATA]) -> 'C:\Users\Victim\AppData\Roaming'
[FILESYSTEM] CopyFileA(src='C:\Users\Public\Astaroth.exe', dst='C:\Users\Victim\AppData\Roaming\5kidRo0t.exe', bFailIfExists=0)
[PERSISTENCE] RegSetValueExA(HKEY_CURRENT_USER\Software\Microsoft\Windows\CurrentVersion\Run, value='5kidRo0t', data='C:\Users\Victim\AppData\Roaming\5kidRo0t.exe')
[PERSISTENCE] RegSetValueExA(HKEY_LOCAL_MACHINE\Software\Microsoft\Windows\CurrentVersion\Run, value='5kidRo0t', data='C:\Users\Victim\AppData\Roaming\5kidRo0t.exe')
[FILESYSTEM] SetFileAttributesA(file='C:\Users\Victim\AppData\Roaming\5kidRo0t.exe', attr=0x6 [HIDDEN | SYSTEM])
\end{lstlisting}
\end{minipage}


\subsection{Costruzione di \texttt{struct hostent} e Scansione ICMP della Sottorete Locale (\texttt{/24})}

\textbf{Funzionamento della Tecnica:} Nella funzione \texttt{0x00401701}, Astaroth verifica la presenza di una scheda di rete operativa invocando \texttt{gethostname()} e passando il nome ottenuto a \texttt{gethostbyname()}. Se l'ambiente isolato restituisce un puntatore nullo o incompleto, l'istruzione assembly \texttt{mov eax, [eax+0xc]} (che accede al campo \texttt{h\_\allowbreak{}addr\_\allowbreak{}list} della struttura C \texttt{hostent}) provoca il crash immediato del processo. Ottenuto un IP valido (es. \texttt{192.168.1.100}), il malware ne estrae i primi tre ottetti, apre un \textit{Raw Socket} ICMP (\texttt{AF\_INET = 2}, \texttt{SOCK\_\allowbreak{}RAW = 3}, \texttt{IPPROTO\_\allowbreak{}ICMP = 1}) e cicla da \texttt{1} a \texttt{254} formattando ciascun indirizzo tramite un wrapper MinGW di \texttt{sprintf} (\texttt{0x004014e0}) e inviando 254 pacchetti \textit{ICMP Echo Request} da 32 byte intervallati da \texttt{Sleep(10)}.

\textbf{Logica di Risoluzione e Implementazione:} Nell'hook di \texttt{gethostbyname} \`e stata allocata sull'heap di Qiling una struttura binaria \texttt{hostent} a 32 bit completa di puntatori incrociati verso l'indirizzo \texttt{192.168.1.100} (Listato~\ref{lst:astaroth_net}). Inoltre, poich\'e la routine interna \texttt{\_\_\allowbreak{}mingw\_\allowbreak{}pformat} richiamata da \texttt{0x004014e0} invocava funzioni di localizzazione DBCS di \texttt{msvcrt.dll} non implementate in Qiling, \`e stato applicato un hook su indirizzo a \texttt{0x004014e0} che legge direttamente dallo stack (\texttt{ESP}) i quattro ottetti, scrive in memoria la stringa ASCII \texttt{"192.168.1.\allowbreak{}X"} e salta all'istruzione di ritorno \texttt{0x0040150c}, consentendo di registrare integralmente la scansione dei 254 host (Listato~\ref{lst:astaroth_net_out}).

\begin{lstlisting}[style=pythoncode, caption={Costruzione manuale in RAM della struttura binaria hostent a 32 bit per sbloccare la scansione ICMP.}, label={lst:astaroth_net}]
@winsdkapi(cc=STDCALL, params={"name": LPCSTR})
def hook_gethostbyname(ql, address, params):
    base = ql.os.heap.alloc(0x80)
    h_name_ptr, h_addr_list_ptr, in_addr_ptr = base + 0x20, base + 0x40, base + 0x50

    ql.mem.write(h_name_ptr, b"DESKTOP-SANDBOX\x00")
    ql.mem.write(in_addr_ptr, py_socket.inet_aton("192.168.1.100"))
    ql.mem.write(h_addr_list_ptr, struct.pack("<II", in_addr_ptr, 0))
    # Layout struct hostent (32-bit): h_name, h_aliases, h_addrtype=2, h_length=4, h_addr_list
    ql.mem.write(base, struct.pack("<IIhhI", h_name_ptr, 0, 2, 4, h_addr_list_ptr))
    return base
\end{lstlisting}


\noindent\begin{minipage}{\linewidth}
\begin{lstlisting}[style=consoleout, caption={Log della risoluzione hostent e del Ping Sweep ICMP sui 254 indirizzi della sottorete locale.}, label={lst:astaroth_net_out}]
[NETWORK] WSAStartup(version=0x202)
[NETWORK] gethostname() -> 'DESKTOP-SANDBOX'
[NETWORK] gethostbyname('DESKTOP-SANDBOX') -> 192.168.1.100 (struct hostent @ 0x500128e)
[NETWORK] socket(af=2 [AF_INET], type=3 [SOCK_RAW], protocol=1 [IPPROTO_ICMP]) -> handle=0x1337
[NETWORK] sendto(sock=0x1337, len=32 [ICMP Echo Request], target=192.168.1.1) (#1/254)
[NETWORK] sendto(sock=0x1337, len=32 [ICMP Echo Request], target=192.168.1.2) (#2/254)
[NETWORK] ... (scansione ICMP ping sweep in corso su 192.168.1.3 .. 192.168.1.253) ...
[NETWORK] sendto(sock=0x1337, len=32 [ICMP Echo Request], target=192.168.1.254) (#254/254)
\end{lstlisting}
\end{minipage}


\subsection{Neutralizzazione dell'Attacco Finale per Esaurimento Risorse (Fork-Bomb e Saturazione RAM)}

\textbf{Funzionamento della Tecnica:} La subroutine finale (\texttt{0x00401912}) coniuga il sabotaggio della macchina vittima con una severa tecnica Anti-Sandbox per esaurimento delle risorse e del tempo di analisi:
\begin{enumerate}
	\item avvia un ciclo infinito che invoca ripetutamente \texttt{Create\allowbreak{}ProcessA} sul proprio eseguibile (\textit{Fork-Bomb}), generando una crescita esponenziale di processi;
	\item ad ogni iterazione riuscita, invoca \texttt{malloc\allowbreak{}(0x17d78400)} allocando \textbf{400.000.000 di byte} ($\approx 381\textrm{ MB}$) di memoria heap e, all'indirizzo \texttt{0x004019d2}, esegue un ciclo di \textbf{100.000.000 di iterazioni} (\texttt{cmpl \$0x5f5e0ff, -0xc(\%ebp)}) che scrive un intero a 32 bit in ciascuna cella allocata. Questa scrittura forzata impedisce l'allocazione \textit{lazy} del sistema operativo, obbligando le VM ad allocare fisicamente le pagine fino al crash per \textit{Out-Of-Memory} e costringendo gli emulatori a tradurre centinaia di milioni di istruzioni fino al \textit{timeout};
	\item inserisce inoltre ritardi \texttt{Sleep(500)} a ogni iterazione e \texttt{Sleep(10000)} in uscita.
\end{enumerate}

\textbf{Logica di Risoluzione e Implementazione:} Per documentare sperimentalmente la natura iterativa dell'attacco senza esaurire la RAM dell'host n\'e bloccare la CPU emulata, l'harness combina quattro difese coordinate (Listato~\ref{lst:astaroth_dos}): ridimensiona trasparentemente in \texttt{hook\_malloc} ogni allocazione $\ge 256\textrm{ MB}$ a un buffer di sicurezza di soli \texttt{4 KB} (\texttt{0x1000}); intercetta l'inizio del ciclo di riempimento a \texttt{0x004019d2} spostando direttamente \texttt{EIP} all'istruzione \texttt{0x004019fc} immediatamente successiva al loop; consente le prime due iterazioni di \texttt{Create\allowbreak{}ProcessA} (restituendo \texttt{1}) per registrare la telemetria della Fork-Bomb e restituisce \texttt{0} alla terza chiamata per forzare l'uscita pulita dal ciclo \texttt{while}; infine azzera tutte le pause \texttt{Sleep} (Listato~\ref{lst:astaroth_dos_out}).

\begin{lstlisting}[style=pythoncode, caption={Contenimento di malloc(381 MB), salto del loop da 100M di scritture e arresto controllato della Fork-Bomb.}, label={lst:astaroth_dos}]
@winsdkapi(cc=CDECL, params={"size": DWORD})
def hook_malloc(ql, address, params):
    size = params["size"]
    if size >= 0x10000000:  # Intercetta malloc(0x17d78400) -> ~381 MB
        return ql.os.heap.alloc(0x1000)  # Buffer ridotto di sicurezza (4 KB)
    return ql.os.heap.alloc(size if size > 0 else 16)

def hook_bypass_ram_fill_loop(ql):
    # Salta le 100.000.000 di scritture int32 da 0x004019D2 direttamente a 0x004019FC
    ql.arch.regs.eip = 0x004019FC

@winsdkapi(cc=STDCALL, params={
    "lpApplicationName": LPCSTR, "lpCommandLine": LPCSTR, "lpProcessAttributes": POINTER,
    "lpThreadAttributes": POINTER, "bInheritHandles": INT, "dwCreationFlags": DWORD,
    "lpEnvironment": POINTER, "lpCurrentDirectory": LPCSTR, "lpStartupInfo": POINTER,
    "lpProcessInformation": POINTER
})
def hook_CreateProcessA(ql, address, params):
    global create_process_count
    create_process_count += 1
    return 1 if create_process_count <= 2 else 0  # Alla 3a iterazione esce dal loop infinito
\end{lstlisting}


\noindent\begin{minipage}{\linewidth}
\begin{lstlisting}[style=consoleout, caption={Esecuzione controllata delle iterazioni Fork-Bomb/RAM-Exhaustion e completamento naturale del processo.}, label={lst:astaroth_dos_out}]
[DOS / FORK-BOMB] CreateProcessA(app='C:\Users\Public\Astaroth.exe') [Iterazione Fork-Bomb #1] -> Consentito (Simulato)
[DOS / RAM-EXHAUSTION] malloc(0x17d78400 [381 MB]) intercettato! Allocato buffer ridotto di sicurezza per prevenire OOM.
[DOS / RAM-EXHAUSTION] Rilevato ciclo di saturazione RAM a 0x4019d2 (100.000.000 scritture int32 = 400 MB). Bypass a 0x4019fc!
[ANTI-DELAY] Sleep(500 ms) intercettato e azzerato
[DOS / FORK-BOMB] CreateProcessA(app='C:\Users\Public\Astaroth.exe') [Iterazione Fork-Bomb #2] -> Consentito (Simulato)
[DOS / RAM-EXHAUSTION] malloc(0x17d78400 [381 MB]) intercettato! Allocato buffer ridotto di sicurezza per prevenire OOM.
[DOS / RAM-EXHAUSTION] Rilevato ciclo di saturazione RAM a 0x4019d2 (100.000.000 scritture int32 = 400 MB). Bypass a 0x4019fc!
[DOS / FORK-BOMB] CreateProcessA(app='C:\Users\Public\Astaroth.exe') [Iterazione #3] -> Interrotto (Restituito 0 per uscire dal loop infinito)
[ANTI-DELAY] Sleep(10000 ms) intercettato e azzerato
[+] EMULAZIONE COMPLETATA CON SUCCESSO FINO ALLA TERMINAZIONE DEL PROCESSO!
\end{lstlisting}
\end{minipage}


\subsection{Artefatti Estratti e Report Telemetrico JSON di Astaroth}

Al termine dell'analisi (completata in pochi secondi effettivi), lo script ha esportato su disco il binario PE decompattato \texttt{astaroth\_\allowbreak{}unpacked\_\allowbreak{}dump.\allowbreak{}bin} (\texttt{86.016 byte}) e il riepilogo telemetrico strutturato \texttt{astaroth\_\allowbreak{}report.\allowbreak{}json} (\texttt{2.744 byte}), riportato integralmente nel Listato~\ref{lst:astaroth_json}.

\begin{lstlisting}[style=consoleout, caption={Report telemetrico JSON generato al termine dell'emulazione di Astaroth (astaroth\_report.json).}, label={lst:astaroth_json}]
{
    "sample": "Astaroth.exe",
    "packer": "UPX 4.22 (PE32 i386)",
    "unpacked_oep": "0x004014c0 / main @ 0x00401a18",
    "decrypted_strings": ["Molon Labe"],
    "crypto_operations": [
        "CryptAcquireContextA(dwProvType=0x18 [PROV_RSA_AES], dwFlags=0xf0000000) -> hProv=0xcafe0001",
        "CryptGenKey(Algid=0x6610 [CALG_AES_256], dwFlags=0x1) -> hKey=0xcafe0002",
        "CryptEncrypt(pbData=0xffffcf7d, len=11, bufLen=27, plaintext='Molon Labe')",
        "CryptDestroyKey(hKey=0xcafe0002)",
        "CryptReleaseContext(hProv=0xcafe0001)"
    ],
    "anti_debug": ["IsDebuggerPresent() -> 0"],
    "privilege_escalation": [
        "ShellExecuteA(op='runas', file='Astaroth.exe', nShowCmd=1) [Tentativo di elevazione privilegi UAC]"
    ],
    "filesystem_actions": [
        "CopyFileA(src='C:\\Users\\Public\\Astaroth.exe', dst='C:\\Users\\Victim\\AppData\\Roaming\\5kidRo0t.exe', bFailIfExists=0)",
        "SetFileAttributesA(file='C:\\Users\\Victim\\AppData\\Roaming\\5kidRo0t.exe', attr=0x6 [HIDDEN | SYSTEM])"
    ],
    "registry_persistence": [
        "RegSetValueExA(HKEY_CURRENT_USER\\Software\\Microsoft\\Windows\\CurrentVersion\\Run, value='5kidRo0t', data='C:\\Users\\Victim\\AppData\\Roaming\\5kidRo0t.exe')",
        "RegSetValueExA(HKEY_LOCAL_MACHINE\\Software\\Microsoft\\Windows\\CurrentVersion\\Run, value='5kidRo0t', data='C:\\Users\\Victim\\AppData\\Roaming\\5kidRo0t.exe')"
    ],
    "network_discovery": {
        "hostname": "DESKTOP-SANDBOX",
        "local_ip": "192.168.1.100",
        "socket_type": "socket(af=2 [AF_INET], type=3 [SOCK_RAW], protocol=1 [IPPROTO_ICMP])",
        "icmp_sweep_targets": ["192.168.1.1", "192.168.1.2", "192.168.1.3", "192.168.1.4", "192.168.1.5", "192.168.1.254"],
        "total_icmp_packets": 254
    }
}
\end{lstlisting}

% =====================================================================
% SEZIONE 8
% =====================================================================
\section{Caso di Studio 4: Analisi di Latrodectus e Studio dei Limiti Strutturali dell'Emulazione User-Mode}

Il quarto e ultimo caso di studio riguarda \textbf{Latrodectus}, un moderno \textit{loader} multi-stadio a 64 bit selezionato specificamente per documentare e analizzare sul piano scientifico i \textbf{casi con esito negativo} nell'emulazione dinamica puramente \textit{user-mode}. In un lavoro sperimentale rigoroso, l'indagine sulle cause esatte per cui determinate contromisure evasive resistono a un emulatore di processo come Qiling risulta altrettanto preziosa dei successi ottenuti sui primi tre campioni, poich\'e consente di tracciare il confine netto tra ci\`o che \`e risolvibile tramite \textit{API/instruction hooking} in spazio utente e ci\`o che richiede necessariamente l'intervento di un vero kernel Windows o l'integrazione con tecniche di estrazione statica.

\subsection[Mascheramento PE32+, Falso DllMain ed Entry Point Esportato]{Mascheramento PE32+, Falso \texttt{DllMain} e Avvio Forzato (\texttt{0x180065709})}

\textbf{Funzionamento della Tecnica:} Pur essendo distribuito con estensione \texttt{.exe}, l'esame dell'intestazione PE32+ a 64 bit rivela che il file \`e in realt\`a una libreria a collegamento dinamico (\texttt{IMAGE\_\allowbreak{}FILE\_\allowbreak{}DLL} attivo nel campo \texttt{Characteristics}). Per trarre in inganno le sandbox automatizzate che caricano una DLL invocandone unicamente il punto d'ingresso canonico (\texttt{DllMain}), gli autori vi hanno collocato una funzione minimale priva di logica operativa che restituisce immediatamente \texttt{1} (\texttt{TRUE}) e termina senza sollevare alcun sospetto. L'ispezione statica della \textit{Export Directory} ha svelato che la vera routine di innesco del primo stadio \`e celata dietro ben \textbf{203 funzioni esportate con nomi casuali}, tutte convergenti sul medesimo indirizzo virtuale interno: \texttt{0x180065709}.

\textbf{Configurazione dell'Harness Sperimentale:} Nello script \texttt{latrodectus\_\allowbreak{}stage1\_\allowbreak{}emulation\_\allowbreak{}harness.\allowbreak{}py} sono stati predisposti gli hook di monitoraggio sul caricamento dinamico dei moduli (\texttt{Load\allowbreak{}LibraryA/\allowbreak{}W}, \texttt{Get\allowbreak{}Proc\allowbreak{}Address}, \texttt{Create\allowbreak{}FileA/\allowbreak{}W}) e i presidi anti-stallo temporale (\texttt{Sleep}, \texttt{GetTickCount}), scavalcando il falso \texttt{DllMain} mediante l'avvio forzato dell'emulazione direttamente all'indirizzo \texttt{STAGE1\_\allowbreak{}EXPORT\_\allowbreak{}VA = 0x180065709} (Listato~\ref{lst:latrodectus_harness}).

\begin{lstlisting}[style=pythoncode, caption={Forzatura dell'esecuzione oltre il falso DllMain verso il vero entry point esportato (0x180065709).}, label={lst:latrodectus_harness}]
STAGE1_EXPORT_VA = 0x180065709

def setup_latrodectus_hooks(ql):
    ql.os.set_api("LoadLibraryA", hook_LoadLibraryA)
    ql.os.set_api("LoadLibraryW", hook_LoadLibraryW)
    ql.os.set_api("GetProcAddress", hook_GetProcAddress)
    ql.os.set_api("CreateFileA", hook_CreateFileA)
    ql.os.set_api("CreateFileW", hook_CreateFileW)
    ql.os.set_api("Sleep", hook_Sleep)
    ql.os.set_api("GetTickCount", hook_GetTickCount)

def avvia_emulazione_base(pe_path, rootfs_path, timeout_sec=10):
    ql = Qiling([pe_path], rootfs_path, verbose=QL_VERBOSE.DEBUG)
    setup_latrodectus_hooks(ql)
    timeout_us = int(timeout_sec * 1_000_000) if timeout_sec > 0 else 0
    print(f"[*] Bypass DllMain vuoto -> Esecuzione forzata a: {hex(STAGE1_EXPORT_VA)}")
    ql.run(begin=STAGE1_EXPORT_VA, timeout=timeout_us)
\end{lstlisting}

\subsection{Offuscamento Vettoriale SIMD, Ciclo Anti-Emulazione (>456M Iterazioni) ed Estrazione Statica}

\textbf{Causa del Primo Blocco Dinamico:} Avviando l'emulazione a \texttt{0x180065709}, nessuno degli hook API registrati viene mai raggiunto prima dello scadere del timeout (Listato~\ref{lst:latrodectus_out}). L'analisi disassemblata del codice macchina ha chiarito la natura squisitamente anti-emulazione dell'ostacolo: a differenza di PlugX, del RAT in Go generico e di Astaroth --- che per attendere o interrogare il sistema invocavano API di Windows intercettabili --- il primo stadio di Latrodectus affoga ogni singola istruzione utile in una densa sequenza di calcoli vettoriali SIMD inutili (\textit{junk instructions}) ed esegue immediatamente all'ingresso un \textbf{doppio ciclo matematico annidato di oltre 456 milioni di iterazioni} puramente aritmetiche sui registri interni della CPU, senza effettuare alcuna chiamata di sistema. Su un processore fisico moderno dotato di pipeline superscalari ed esecuzione \textit{out-of-order}, tale ciclo si esaurisce in una frazione di secondo; al contrario, in un motore di traduzione binaria software come Unicorn/Qiling, ogni singola iterazione deve essere decodificata ed emulata via software, causando inevitabilmente lo stallo per superamento del tempo massimo. Ulteriori test condotti forzando manualmente il registro \texttt{RIP} oltre il ciclo matematico hanno evidenziato una seconda trappola strutturale: il codice ispeziona la lista concatenata dei moduli caricati in memoria partendo dalla propria intestazione PE, all'interno della quale gli autori hanno inserito intenzionalmente un puntatore corrotto che il loader nativo del kernel Windows tollera, ma che provoca un'eccezione fatale nel loader user-mode di Qiling.


\noindent\begin{minipage}{\linewidth}
\begin{lstlisting}[style=consoleout, caption={Evidenza sperimentale dello stallo computazionale anti-emulazione sul primo stadio di Latrodectus.}, label={lst:latrodectus_out}]
=================================================================
[*] Avvio Analisi Dinamica: Latrodectus.exe
[*] Target SHA-256 atteso : fc58ae120d35b751525f6c312f9aeff75e9151f02fa113815bd01acc49f63596
=================================================================
[*] Bypass DllMain vuoto -> Esecuzione forzata a: 0x180065709
[*] Esecuzione in corso (timeout impostato a 10s)...
-----------------------------------------------------------------
[*] Esecuzione terminata per timeout (nessuna API invocata; RIP bloccato nel loop matematico SIMD).
\end{lstlisting}
\end{minipage}


\textbf{Recupero del Secondo Stadio tramite Analisi ed Estrazione Statica:} Isolata staticamente nel disassemblato la routine posta a valle del ciclo di rallentamento, si \`e accertato che il primo stadio localizza nella sezione delle risorse (\texttt{.rsrc}) del file un blob cifrato di \textbf{102.914 byte} e lo decifra in memoria mediante un algoritmo XOR ciclico con chiave di 20 byte. Replicando esattamente tale algoritmo in uno script Python di decifratura statica, \`e stato estratto integralmente e senza corruzioni l'eseguibile del secondo stadio.

\subsection{Controllo del Processo Ospite, Lettura di \texttt{ntdll.dll} da Disco e \textit{Direct Syscalls}}

L'ispezione tecnica del secondo stadio cos\`i recuperato ha portato alla luce i due meccanismi architetturali che rendono strutturalmente impossibile proseguire l'emulazione dinamica in Qiling:

\begin{enumerate}
	\item \textbf{Validazione rigorosa del processo contenitore:} All'avvio, il secondo stadio controlla il nome del processo ospite in cui \`e mappato, proseguendo unicamente se rileva i binari di sistema legittimi di Windows \texttt{rundll32.exe} o \texttt{regsvr32.exe} e terminando istantaneamente in qualsiasi altro contesto.
	\item \textbf{Rimozione totale dei ganci (\textit{Unhooking}) e \textit{Direct Syscalls} lette da disco:} Consapevoli del fatto che sandbox, soluzioni EDR ed emulatori user-mode (incluso Qiling) intercettano l'attivit\`a dei programmi agganciando le funzioni esportate dalle copie delle librerie di sistema gi\`a caricate in RAM (in primo luogo \texttt{ntdll.dll}), gli sviluppatori di Latrodectus ignorano deliberatamente la copia di \texttt{ntdll.dll} presente nello spazio di indirizzamento del processo. Il secondo stadio apre invece in lettura diretta dal file system il binario pulito \texttt{C:\textbackslash\allowbreak{}Windows\textbackslash\allowbreak{}System32\textbackslash\allowbreak{}ntdll.\allowbreak{}dll}, ne riversa il contenuto grezzo in un buffer di memoria privato allocato dinamicamente, individua al suo interno gli \textit{stub} binari delle primitive native di gestione delle sezioni (\texttt{Nt\allowbreak{}Create\allowbreak{}Section}, \texttt{Nt\allowbreak{}Map\allowbreak{}View\allowbreak{}Of\allowbreak{}Section}) ed esegue direttamente l'istruzione assembly \texttt{syscall} dall'interno di quel buffer privato per mappare il terzo stadio.
\end{enumerate}

Questa architettura mette a nudo il limite invalicabile di qualunque emulatore puramente \textit{user-mode}: Qiling intercetta le invocazioni alle API agganciando gli indirizzi esportati dalle DLL mappate dal proprio loader, ma \textbf{non dispone di un vero kernel Windows a Ring 0} n\'e di una \textit{System Service Descriptor Table} (SSDT) completa in grado di servire transizioni \texttt{syscall} arbitrarie originate da buffer privati costruiti leggendo i file di sistema dal disco. Quando la CPU emulata incontra l'istruzione \texttt{syscall} nel buffer privato di Latrodectus, tutti gli hook Python vengono completamente scavalcati e l'emulazione dinamica non pu\`o che arrestarsi.

\subsection{Recupero Statico del Terzo Stadio (\texttt{UpdaterTag.dll}) e Conclusioni sul Caso Negativo}

Proseguendo con metodiche di estrazione statica sul secondo stadio \`e stato infine recuperato anche il payload di terzo stadio, identificato internamente come \texttt{Updater\allowbreak{}Tag.\allowbreak{}dll}. Esso presenta un offuscamento estremo: \`e privo di qualsiasi stringa leggibile su disco (ogni stringa viene decifrata sullo stack esclusivamente nell'istante in cui la singola funzione viene richiamata) e risolve tutte le funzioni di Windows a runtime tramite codici numerici precalcolati (\textit{API hashing}), azzerando il contenuto informativo della tabella delle importazioni.

Nel suo insieme, Latrodectus opera dunque come un caricatore \textit{fileless} a tre stadi in cui nessun modulo intermedio o finale tocca mai il disco in chiaro: il primo stadio manda in timeout gli emulatori software tramite saturazione algebrica e vettoriale della CPU; il secondo neutralizza qualsiasi \textit{hooking} user-mode rileggendo \texttt{ntdll.dll} da disco ed eseguendo \textit{direct syscalls}; il terzo (\texttt{Updater\allowbreak{}Tag.\allowbreak{}dll}) esfiltra informazioni dall'host e contatta l'infrastruttura C2 per il rilascio di ulteriori minacce. Questo caso negativo dimostra concretamente che, di fronte a offuscamento computazionale pesante combinato con \textit{direct syscalls}, l'emulazione user-mode va affiancata a decifratura statica mirata per lo spacchettamento degli involucri e a macchine virtuali Windows complete dotate di monitoraggio a livello del kernel o dell'hypervisor (\textit{Virtual Machine Introspection}) per l'osservazione comportamentale dell'ultimo stadio.

% =====================================================================
% SEZIONE 9
% =====================================================================
\section{Analisi Comparativa dei Risultati, Valutazione Metodologica e Conclusioni}

\subsection{Matrice Comparativa delle Tecniche Anti-VM/Anti-Sandbox e delle Contromisure}

La Tabella~\ref{tab:comparative_matrix} offre una sintesi strutturata di tutte le tecniche di evasione, anti-emulazione e anti-analisi affrontate nei quattro casi di studio, mettendo a confronto il meccanismo offensivo impiegato dal malware con la specifica strategia ingegnerizzata negli script Qiling/Python.

\begin{table}[!htbp]
	\centering
	\scriptsize
	\setlength{\tabcolsep}{4.5pt}
	\renewcommand{\arraystretch}{1.32}
	\begin{tabularx}{\textwidth}{|P{1.85cm}|P{2.55cm}|Y|Y|}
		\hline
		\textbf{Campione} & \textbf{Categoria Difesa} & \textbf{Meccanismo Tecnico del Malware} & \textbf{Strategia di Risoluzione Implementata} \\
		\hline
		\textbf{PlugX} & Evasione temporale & Confronto \texttt{GetTickCount} prima/dopo \texttt{Sleep(5000)} con loop infinito & Azzeramento attesa in \texttt{hook\_\allowbreak{}Sleep} e somma differita dei millisecondi in \texttt{hook\_\allowbreak{}GetTickCount} \\
		\hline
		\textbf{PlugX} & Interazione umana (HID) & Doppio campionamento \texttt{GetCursorPos} e verifica \texttt{GetForeground\allowbreak{}Window} & Iniezione jitter $[-15, +15]$ px in RAM e fornitura di \texttt{HWND} crescenti sfruttando il salto incondizionato \\
		\hline
		\textbf{PlugX} & Pre-infezione / C I/O & Sondaggio \texttt{fopen("rb")} su \texttt{MediaCenter.\allowbreak{}exe} e scrittura via \texttt{fwrite} & Custom VFS in Python (\texttt{fopen -> 0} in lettura; cattura automatica dei 74.358 byte su \texttt{fwrite}) \\
		\hline
		\textbf{PlugX} & Persistenza post-reboot & Controllo percorso di lancio tramite \texttt{GetModule\allowbreak{}FileNameA} & Scrittura in memoria del path atteso \texttt{MicroMedia\textbackslash\allowbreak{}MediaCenter.exe} e cattura C2 HTTP (IE8) \\
		\hline
		\textbf{RAT in Go generico} & Spawn multi-processo & Verifica \texttt{\_HIDDEN\_\allowbreak{}CHILD}, lancio copia con \texttt{Create\allowbreak{}ProcessW} e \texttt{os.Exit(0)} & Hook \texttt{Create\allowbreak{}ProcessW} + azzeramento chirurgico \texttt{RAX=0} a \texttt{0x77dd94} per transizione in-process \\
		\hline
		\textbf{RAT in Go generico} & Fingerprinting WMI & Query COM/OLE \texttt{SELECT * FROM Win32\_\allowbreak{}Processor} & Costruzione VTable \texttt{IDispatch}/\allowbreak{}\texttt{IEnumVARIANT} a 64 bit in RAM (spoofing Intel Core i7-9700K 8 core) \\
		\hline
		\textbf{RAT in Go generico} & Tempo / Carico CPU Go & Lettura \texttt{Get\allowbreak{}Tick\allowbreak{}Count64}, \texttt{KUSER\_\allowbreak{}SHARED\_\allowbreak{}DATA} (\texttt{0x7ffe0008}) e delta CPU 1 s & Avanzamento programmatico \texttt{KUSER\_\allowbreak{}SHARED\_\allowbreak{}DATA}, uptime 3604 s e carico simulato al 24,69\% \\
		\hline
		\textbf{RAT in Go generico} & Risorse e Rete / C2 & Controllo RAM/Disco/NIC, IP/GeoIP, scan 23 porte e C2 WebSocket Protobuf & Profilo 16 GB RAM / 500 GB NTFS / 154 MB RX, mock HTTP GeoIP e server C2 RFC 6455 + Protobuf \\
		\hline
		\textbf{Astaroth} & Packing e CryptoAPI & Compressione UPX 4.22, XOR \texttt{0x5A} e cifratura \texttt{AES-256} (\texttt{cryptsp.dll}) & Unpacking naturale fino a \texttt{0x00401a18} (dump PE) e 5 hook CryptoAPI con cattura di \texttt{"Molon Labe"} \\
		\hline
		\textbf{Astaroth} & Anti-Debug / UAC / Rete & \texttt{IsDebugger\allowbreak{}Present}, \texttt{Shell\allowbreak{}ExecuteA\allowbreak{}("runas")} e ping sweep \texttt{/24} & Ritorno \texttt{0} e \texttt{42}, costruzione manuale \texttt{struct hostent} in heap e bypass \texttt{sprintf} MinGW \\
		\hline
		\textbf{Astaroth} & Sabotaggio / DoS & Loop infinito \texttt{Create\allowbreak{}ProcessA} + \texttt{malloc(381 MB)} + 100M scritture int32 & Riduzione \texttt{malloc} a 4 KB, salto \texttt{EIP} (\texttt{0x4019d2 -> 0x4019fc}) e uscita alla 3a iterazione \\
		\hline
		\textbf{Latrodectus} & Anti-Emulazione / Kernel & Falso \texttt{DllMain}, loop SIMD $>456\textrm{M}$ cicli, lettura raw \texttt{ntdll.dll} e \texttt{syscall} & Avvio forzato a \texttt{0x180065709}, documentazione limite user-mode ed estrazione statica XOR Stadi 2--3 \\
		\hline
	\end{tabularx}
	\caption{Matrice comparativa delle difese analizzate e delle soluzioni sviluppate nei quattro casi di studio.}
	\label{tab:comparative_matrix}
\end{table}

\subsection{Conclusioni Generali dello Studio}

Il percorso sperimentale condotto sui campioni \textbf{PlugX}, \textbf{RAT in Go generico}, \textbf{Astaroth} e \textbf{Latrodectus} mediante \textbf{Qiling 1.4.6 (su Python 3.12)} all'interno di un ambiente virtuale isolato (\texttt{env}) corredato dalle librerie DLL di Windows 10 (32/64 bit) e Windows 7 (32 bit) consente di trarre conclusioni metodologiche chiare sulle potenzialit\`a e sui confini dell'emulazione binaria dinamica mirata:

\begin{enumerate}
	\item \textbf{Efficacia chirurgica e velocit\`a sull'evasione basata su API e stato ambientale:} Su \textbf{PlugX}, sul \textbf{RAT in Go generico} e su \textbf{Astaroth}, l'approccio basato su Qiling ha permesso di neutralizzare in pochi secondi controlli che in una sandbox tradizionale avrebbero richiesto lunghe attese, complesse riconfigurazioni dell'hypervisor o l'esposizione di infrastrutture di rete. La possibilit\`a di leggere i buffer in memoria nell'istante esatto della chiamata API (catturando ad esempio i 74.358 byte decifrati di PlugX su \texttt{fwrite}, la stringa effimera \texttt{"Molon Labe"} in Astaroth prima della cifratura AES-256 o i pacchetti binari Protocol Buffers del RAT in Go generico prima della trasmissione in rete) offre una visibilit\`a interna difficilmente eguagliabile da un monitoraggio esterno di tipo \textit{black-box}.
	
	\item \textbf{Capacit\`a di simulare interi ecosistemi runtime e protocolli C2 moderni:} Il caso di studio sul \textbf{RAT in Go generico} ha dimostrato come, pur partendo da un emulatore privo di servizi COM/WMI, di timer kernel su \texttt{KUSER\_\allowbreak{}SHARED\_\allowbreak{}DATA} e di rete esterna, un'attenta ingegnerizzazione in Python (dalle VTable binarie a 64 bit fino a un server interattivo WebSocket RFC 6455 + Protocol Buffers) consenta di portare a compimento l'esecuzione di complessi eseguibili moderni scritti in Go 1.23, collaudandone in totale sicurezza tutte le funzionalit\`a di esfiltrazione e gestione remota dei file.
	
	\item \textbf{Resilienza agli attacchi distruttivi contro l'ambiente di analisi:} L'indagine su \textbf{Astaroth} ha evidenziato un ulteriore punto di forza decisivo dell'emulazione controllata rispetto alle macchine virtuali tradizionali: di fronte a una \textit{Fork-Bomb} combinata con allocazioni ripetute da 381--400 MB e cicli da 100 milioni di scritture progettati per provocare il crash per esaurimento della RAM, il controllo diretto sull'allocatore \texttt{malloc} e sul registro \texttt{EIP} ha permesso di osservare e registrare l'attacco mantenendo il consumo effettivo di memoria a pochi kilobyte.
	
	\item \textbf{Consapevolezza rigorosa dei limiti strutturali dell'user-mode:} Infine, l'eccezione POSIX emersa su \texttt{SysWOW64} con \textit{PECompact2} nello Stadio 2 di PlugX e, soprattutto, l'esito negativo documentato su \textbf{Latrodectus} definiscono con precisione i limiti intrinseci dell'emulazione a livello utente. Quando un campione combina cicli puramente computazionali da centinaia di milioni di istruzioni prive di chiamate API con tecniche di \textit{unhooking} e \textit{direct syscalls} invocate da copie di \texttt{ntdll.dll} lette direttamente dal disco, l'assenza di una CPU fisica dedicata e di un vero kernel Windows a Ring 0 impedisce la prosecuzione dinamica in Qiling. In tali scenari avanzati, la metodologia ottimale risiede in un approccio ibrido: impiegare l'emulazione mirata e gli script di decifratura statica per rimuovere rapidamente gli strati esterni di offuscamento, demandando l'osservazione dinamica del modulo finale a sandbox full-system equipaggiate con introspezione a livello kernel o hypervisor.
\end{enumerate}

\vspace{1.2cm}
\vfill
\noindent\rule{\textwidth}{0.4pt}\\[4pt]
{\scriptsize\color{muted}\textit{Nota di supervisione tecnica:} La stesura della presente relazione tecnica, cos\`i come l'analisi, la revisione e la supervisione dei codici Python di emulazione dinamica associati allo studio, sono state realizzate con il supporto di agenti di intelligenza artificiale basati sul modello \textbf{Google Gemini 4 Carbon}.}

\end{document}
